\documentclass[journal,twoside,web]{ieeecolor}

\include{Config/config_ieee}
\include{Config/config_des}

% abbrs for this article only
\newcommand{\sFourC}{\ensuremath{\text{S}^4\text{C}}}

\markboth{\hskip25pc IEEE TRANSACTIONS AND JOURNALS TEMPLATE}
{Author \MakeLowercase{\textit{et al.}}: Title}
\begin{document}
\title{Synthesis of Cognitive Non-Blocking Supervisors for Partially-Observed Discrete-Event Systems}
\author{Yiwei Zheng, Bohan Cui, and Xiang Yin
\thanks{TODO This paragraph of the first footnote will contain the date on which you submitted your article for review. It will also contain support information, including sponsor and financial support acknowledgment. 
For example, ``This work was supported in part by the U.S. Department of Commerce under Grant 123456.'' }
\thanks{TODO First A. Author is with the National Institute of Standards and 
Technology, Boulder, CO 80305 USA (e-mail: author@boulder.nist.gov). }
\thanks{TODO Second B. Author Jr. was with Rice University, Houston, TX 77005 USA. 
He is now with the Department of Physics, Colorado State University, Fort Collins, CO 80523 USA (e-mail: author@lamar.colostate.edu).}}

\maketitle

\begin{abstract}
This article investigates the problem of synthesizing cognitive non-blocking supervisors for partially observed discrete event systems (DESs).
Cognitive non-blockingness requires that, for each state of the DES, the system can eventually certify, based solely on partial observations, that it has reached a target state.
The synthesis problem is to construct a supervisor under which the closed-loop system satisfies this observation-based progress requirement.
%
We introduce a new structure, named the bipartite cognitive controller (BCC), to capture both the evolution of information states and the microstate-level transition relations within the unobservable reach.
We prove that the event sequences generated by a BCC can be directly projected as the language of the controlled DES, thereby establishing the equivalence between the BCC-induced language and the controlled DES language.
Moreover, we show that cognitive non-blockingness can be guaranteed by requiring the Y-states to be cognitively live and the Z-states to be potentially feasible.
The resulting BCC can then be projected back to a supervisor for cognitive non-blocking synthesis of the original DES.
\end{abstract}

\begin{IEEEkeywords}
Discrete-event systems (DES), partial observation, non-blocking, supervisory control. %maximally permissiveness.
%\underline{https://www.ieee.org/publications/services/thesaurus-access-page.com}
\end{IEEEkeywords}

\section{Introduction}
\label{sec:intro}
\IEEEPARstart{A}{s} modern computation and networking technologies continue to advance, discrete event systems (DESs) have become increasingly important in various aspects of modern life and industrial production \cite{yin2024formal, fokkink2024offline}.
Owing to their capability to model complex event-driven behaviors, DESs have been widely applied in areas such as automated manufacturing systems, communication networks, transportation systems, and robotic coordination.
%
Non-blockingness is one of the most fundamental properties of DESs, since it guarantees that the system can always complete its prescribed tasks without being trapped before task completion \cite{cassandras2021introduction}.
In general, a DES is said to be non-blocking if every reachable state can still evolve to a marked (or target) state through some feasible event sequence.
This ensures that the system can continuously make progress toward task completion rather than entering an irreversible blocking state, for example, a deadlock or a livelock state.

%
% TODO
% 要不要把enforce这个词引入?
The supervisory control framework introduced by Ramadge and Wonham \cite{ramadge1987supervisory} provides a systematic approach to enforcing desired behaviors of DESs by enabling or disabling controllable events.
In the classical non-blocking synthesis problem, the objective is to construct a supervisor such that every behavior permitted by the closed-loop system can be extended to a marked behavior.
%
However, under partial observation, classical non-blockingness does not guarantee observational certifiability of reaching marked states. Even if a marked state is reachable in the underlying system, the supervisor may be unable to establish this based on partial observations.
%
Let us consider the motivating example of a robotic system deployed in a partially observable environment, which is illustrated in Fig.~\ref{fig:dfa_grid_env} (a). 
The robot has full knowledge of the environment model, while its exact location remains uncertain due to insufficient sensing and limited localization accuracy.
Suppose that the robot is allowed to transmit data only in safe regions, which are regarded as the target regions in this task. 
If a safe target location is observationally indistinguishable from an unsafe location, then physical reachability of a safe region is not sufficient for safe data transmission.
Instead, safety must be certifiable from observations.
Therefore, before transmitting data, the robot must ensure that all states consistent with its observation belong to the set of safe target locations.

This mismatch between non-blockingness and observational certifiability motivates us to introduce the notion of \emph{cognitive non-blockingness}, an observation-based progress property requiring that task completion be not only reachable but also certifiable under partial observation.
%
We then study the synthesis of partial-observation supervisors that enforce cognitive non-blockingness on the closed-loop system.
%
Built upon the supervisory control framework, this article further studies the synthesis of partial-observation supervisors that enforce cognitive non-blockingness on the closed-loop system.
Different from classical non-blocking synthesis, the proposed problem must account for the evolution of information states, since task completion has to be certified from observations rather than merely reached by the actual state.

%
% TODO 争取给放到第二页?
\begin{figure}[t!]
    \centering
    \includegraphics[width=\linewidth]{Figs/Grid_Environment_DFA.pdf}
    \caption{Deterministic finite state automaton $G$ constructed from the motivating example of the robotic system deployed in a partially observable environment from Section~\ref{sec:intro}, where the set of marked states $Q_m = \{ q_3, q_5, q_6, q_{18} \}$, the set of observable events $\Sigma_o = \{ o_1, o_2, o_3 \}$, and the set of controllable events $\Sigma_c = \{ c_1, c_2, o_3 \}$.}
    % \textcolor{red}{not good, the feasible transitions in (a) are not fully captured by (b), e.g., 19 $\to$ 13, 19 $\to$ 15..., maybe show them separately or motivate the paper only use words}
    %
    % Thanks for your comment
    %
    \label{fig:dfa_grid_env}
\end{figure}

The main contributions of this article are summarized as follows.
%
\begin{enumerate}[(i)]
    \item We propose a new observation-based property, named \emph{cognitive non-blocking}, for partially observed DESs.
    Compared with classical non-blockingness studied in  \cite{yin2016synthesis, yin2018synthesis}, the proposed notion strengthens the standard reachability-based requirement by incorporating observational certification of task completion.
    Specifically, it requires not only that marked states are eventually reachable, but also that the supervisor can certify, based solely on partial observations, that the current state belongs to the marked set.

    \item We develop a new structural representation, called \emph{bipartite cognitive controller} (BCC), to characterize both the evolution of information states and the microstate-level transitions hidden within each unobservable reach. 
    Existing bipartite-controller-based approaches \cite{yin2016uniform,zheng2026enforcement} describe closed-loop evolution in terms of control decisions and observations. In contrast, the proposed BCC further refines this representation at the microstate level, thereby enabling the synthesis of cognitive non-blocking supervisors.

    \item We further develop a synthesis method for cognitive non-blocking supervisors based on the BCC structure. 
    This result strengthens the framework of \cite{yin2016synthesis} by iteratively pruning Z-states that are not potentially feasible and Y-states that are not cognitively live, while preserving structural consistency, so that the closed-loop system guarantees cognitive non-blockingness for every reachable state.
    
\end{enumerate}

\subsection{Related works}

% 1 partially observable supervisor for the IS-based property
The supervisory control of DESs under partial observation has been extensively studied in the context of information-state-based (IS-based) properties such as safety \cite{ma2022resilient}, diagnosability \cite{li2025supervisory, hu2023asynchronous}, opacity \cite{lafortune2018history, zheng2026enforcement},
detectability \cite{shu2013enforcing}, anonymity and attractability, etc. \cite{yin2024formal, basile2022supervisory}.
These properties require the synthesis of a supervisor that ensures that all executions of the closed-loop system satisfy the given property.
%
In \cite{yin2016uniform}, a uniform approach is developed to synthesize a bipartite controller that embeds all reachable information states and all admissible supervisory control strategies, which can be further used to synthesize the maximally permissive supervisor for specific properties.
In \cite{xie2023optimal}, this bipartite controller is extended to solve the problem of enforcing opacity for both qualitative (or infinite-step opacity) and quantitative cases.
In \cite{zheng2026enforcement}, an extended bipartite controller is developed for interval weighted DESs for the enforcement of weighted opacity, for which the enabling interval and the sequences of control patterns determine the reachability of the system.
%
However, most existing IS-based properties focus on information-state consistency, where system behavior can be characterized without explicitly reasoning about future progress. 
In contrast, non-blockingness is inherently a progress property that depends on the existence of continuations to marked states and therefore requires reasoning beyond current information states. 
As a result, directly pruning inconsistent information states using classical bipartite controllers is insufficient, which motivates the development of new observation-based progress properties.


% 2 non-blocking problems
Non-blockingness is a fundamental progress requirement in DESs, requiring that every generated behavior can be extended to a marked behavior \cite{chen1991dealing}.
%
The verification of non-blockingness has been extensively investigated for different DES frameworks, including modular DESs \cite{tang2024compositional}, adapted systems \cite{thuijsman2024transformational}, and compositional analysis of extended finite-state machines \cite{mohajerani2016framework}.
%
In \cite{tomas2018complexity}, the computational complexity of verifying non-blockingness in modular DESs is rigorously analyzed, showing that the problem is PSPACE-complete in general, while special cases such as DFA with one shared event or a bounded number of DFAs are NP- or NL-complete.
%
In \cite{pilbrow2015algorithm}, the authors develop a set of compositional abstraction rules, including blocked events, to efficiently verify non-blockingness in modular DESs.
%
In \cite{tang2024compositional}, a framework is developed for compositional non-blocking verification of finite automata with prioritized events, introducing conflict-preserving abstractions and shaped quotient automata to efficiently ensure non-blockingness in modular DESs.
%
However, most of these works on non-blocking analysis focus on verification rather than synthesis. In many practical scenarios, it is necessary to synthesize a supervisor to actively control the system such that its output behaviors are guaranteed to satisfy the non-blocking property \cite{chen1991dealing, malik2013compositional}. 
This motivates the development of synthesis methods that ensure non-blockingness under various system architectures and observation constraints, bridging the gap between theoretical verification results and practical supervisory control design.


% 3 A + B 
In recent research, the combination of non-blockingness and partial-observation supervisory control has also been widely investigated \cite{fokkink2024offline, hou2023distributed, zhang2024quantitatively}.
%
In \cite{yin2016synthesis}, the IS-based bipartite controller of \cite{yin2016uniform} is extended to synthesize the safe, non-blocking, and maximally permissive supervisor.
%
In \cite{xiang2024existence}, a game-theoretic framework is developed to synthesize bounded non-blocking supervisors for DESs by transforming the supervisor existence problem into a reachability game over information states and admissible control decisions.
%
In \cite{jan2024supervisory}, modular supervisory control under partial observation is investigated from the perspective of normality, where a global non-blocking and maximally permissive normal supervisor is achieved through the parallel composition of local normal supervisors.
%
In \cite{zhang2025bounded},  a new bounded-time non-blocking supervisory control problem is solved for timed DES, in which each task must be completed within a bounded time counted by the number of ticks.
%
Although significant progress has been made in integrating non-blockingness with partial-observation supervisory control \cite{fokkink2024offline, hou2023distributed, zhang2024quantitatively}, non-blockingness only guarantees that the system can eventually reach a marked state corresponding to task completion. 
This reveals a fundamental gap between reachability under supervisory control and its certifiability under partial observation, which motivates the introduction of observation-based non-blocking properties that explicitly connect system observations to task progress.


\subsection{Organization}

The remainder of this article is organized as follows.
%
Section~\ref{sec:prelimi} presents the system model and the notions for partial-observation supervisory control.
Section~\ref{sec:prob} introduces the notion of cognitive non-blockingness, gives the problem formulation, and illustrates the challenges of the proposed problem.
Section~\ref{sec:bcc} puts forward the new structure of the  bipartite cognitive controller and establishes its basic properties.
Section~\ref{sec:solu} presents the method for synthesizing supervisors that enforce cognitive non-blockingness.
Section~\ref{sec:special_z_marked} discusses the special case in which the marked states are within the unobservable reach.
Section~\ref{sec:con} concludes the article and gives directions for future research.


\section{Preliminaries}\label{sec:prelimi}

This section introduces the preliminaries for the verification and synthesis of cognitive non-blocking in discrete-event systems (DESs), including system models and supervisory control under partial observation.


\subsection{System model}

We use $\mathbb{N}$, $\mathbb{N}_{\geq 0}$ to denote the sets of natural numbers and nonnegative natural numbers, respectively. 
We use $\varepsilon$ to denote the empty string of an automaton. 
%
The DES considered is modeled as a \emph{deterministic} finite-state automaton (DFA) and formally defined as follows.
\begin{definition}
    A DFA is a five-tuple
    $$
        G = (Q, \Sigma, \delta, q_0, Q_m)
    $$
    where
    $Q$ is the finite set of states, 
    $\Sigma$ is the finite set of events, 
    $\delta \subseteq Q \times \Sigma \times Q$ is the set of transition relation,
    $q_0$ is the initial state, 
    and $Q_m$ is the set of marked states.
    \label{def:DFA}
\end{definition}
%
The \emph{run} of DFA $G$ is a sequence of states and transitions
$
\rho = q_0 \xrightarrow[]{\sigma_0} q_1 \xrightarrow[]{\sigma_1} \cdots q_n \xrightarrow[]{\sigma_n} q_{n+1},
$
where $(q_{i}, \sigma_i, q_{i+1}) \in \delta$ for $0 \leq i \leq n, n \in \mathbb{N}_{\geq 0}$.
The \emph{sequence} corresponding to run $\rho$ is denoted as
$tr(\rho) = \sigma_0 \, \sigma_1 \, \cdots \, \sigma_n \in \Sigma^*$,
%
With a slight abuse of notation, we use $\delta(q, r) = q'$ as the extended partial transition function to denote the reachable state $q'$ from state $q$ and sequence $r$.
%
The behavior generated by $G$ is described by 
$\mathcal{L}(G) = \{ s \in \Sigma^* | \delta(q_0, s)! \}$, 
where $!$ means ``is defined'', and the marked behavior is 
$\mathcal{L}_m(G) = \{ s \in \Sigma^* | \delta(q_0, s) \in Q_m \}$. 
% The prefix-closure of a language $\mathcal{L}(G)$ is
% $$
% \overline{\mathcal{L}(G)} = \{ s \in \Sigma^* | \exists t \in \Sigma^* : st \in \mathcal{L}(G) \},
% $$
% and we say that $\mathcal{L}$ is prefix-closed if $\overline{\mathcal{L}} = \mathcal{L}$.

In DFA $G$, a state $q \in Q$ is said to be \emph{accessible} if there exists a sequence $s \in \mathcal{L}(G)$ such that $\delta(q_0, s) = q$. 
A state $q \in Q$ is \emph{co-accessible} if there exists a sequence $t \in \mathcal{L}(G)$ such that $\delta(q, t) \in Q_m$.
%
A standard progress property in DES is \emph{non-blocking}, requiring that every reachable state can eventually reach a marked state, as formally given in the following definition.
\begin{definition}
    Given a DFA $G=(Q,\Sigma,\delta,q_0,Q_m)$, we say that $G$ is
    \emph{non-blocking} if every accessible state of $G$ is co-accessible to  
    $Q_m$,
    \begin{equation} \label{eq:nb_1}
        \forall s \in \mathcal{L}(G), \exists t \in \Sigma^* \ s.t. \ st \in \mathcal{L}_m(G).
    \end{equation}
    \label{def:nb}
\end{definition}
%
Based on Definition \ref{def:nb}, one can easily check that the system is non-blocking if and only if 
\begin{flalign}
    \forall q\in Q, \
    &(\exists s \in \mathcal{L}(G): \delta(q_0,s)=q)
    \Rightarrow  \notag \\
    &(\exists t\in\Sigma^*: \delta(q,t)\in Q_m).
\label{eq:nb_2}
\end{flalign}
% Intuitively, the notion of non-blockingness means that the system can always continue from any reachable state to complete its task. 
% In other words, no reachable behavior is trapped before reaching a marked state.

We assume the DFA $G$ is partially observed, and the set of events $\Sigma$ can be partitioned into two disjoint subsets, 
$
\Sigma = \Sigma_o \ \dot{\cup} \ \Sigma_{uo},
$ 
where $\Sigma_o$ is the set of observable events, and $\Sigma_{uo}$ is the set of unobservable events.
%
The \emph{natural projection}, $P : \Sigma^* \rightarrow \Sigma_{o}^*$ is recursively defined as
\begin{equation}
P(s \sigma) =
\begin{cases}
P(s), & \text{if } \sigma \in \Sigma_{uo} \\
P(s)\sigma, & \text{otherwise}
\end{cases}
\label{eq:proj}
\end{equation}
with
$P(\varepsilon)  = \varepsilon$.
Intuitively, the projection $P(\cdot)$ prunes the unobservable events within the given sequence $s$ and outputs the sequences of observable events.
With a slight abuse of notation, we denote $P(\mathcal{L}(G))$ as the observable language of DFA $G$.

\subsection{Supervisory control under partial observations}

In the supervisory control framework initiated in \cite{ramadge1987supervisory}, a supervisor is imposed on $G$ to achieve some specification by dynamically enabling or disabling events. 
The event set $\Sigma$ is partitioned into two disjoint subsets: the subset of controllable events $\Sigma_c$, and the subset of uncontrollable events $\Sigma_{uc}$, i.e., 
$\Sigma = \Sigma_c \ \dot{\cup} \ \Sigma_{uc}$.
%
The full set of control patterns for DFA $G$ is defined as $\Gamma = 2^{\Sigma_c}$, where each $\gamma \in \Gamma$ specifies the set of enabled controllable events.
%
For DFA $G$, a partial-observation supervisor $S$ is defined as a function $S : P(\mathcal{L}(G)) \rightarrow \Gamma$, such that for $s, t \in \Sigma^*$, we have $S(P(s))=S(P(t))$ if $P(s) = P(t)$.
%
We use $S / G$ to denote controlled system, and the language generated by $S /G$, denoted by $\mathcal{L}(S / G)$, is recursively defined as follows,
\begin{enumerate}[(i)]
    \item $\varepsilon \in \mathcal{L}(S / G)$;
    \item $s \, \sigma \in \mathcal{L}(S / G)$ if and only if
    $s \in \mathcal{L}(S / G)$ $\wedge$ $s \, \sigma \in \mathcal{L}(G)$ $\wedge$ $\sigma \in S(P(s))$.
\end{enumerate}
We define the marked language of the closed-loop system $S / G$, as $\mathcal{L}_m(S / G) = \mathcal{L}(S / G) \cap \mathcal{L}_m(G)$.
%
It is worth noting that the desired closed-loop property can be achieved by a trivial supervisor that disables controllable events whenever possible, but such a solution is overly conservative and sacrifices behavioral permissiveness. Therefore, we seek supervisors that satisfy the specification while preserving maximal admissible behavior.
%
A supervisor $S$ is said to be maximally permissive within a prescribed class of admissible supervisors if there does not exist another admissible supervisor $S'$ such that
\begin{equation}\label{eq:mps}
    \mathcal{L}(S/G) \not \subset \mathcal{L}(S'/G).
\end{equation}
%
For a sequence $s\in\mathcal{L}(S/G)$, the \emph{state estimate} upon the occurrence of $s \in \mathcal{L}(G)$, namely, the set of states in which the system could be after observing $P(s) \in P(\mathcal{L}(G))$, is defined as
\begin{flalign}\label{eq:esti:1}
\mathcal{E}_S(s) := \{ &q \in Q \mid \exists t \in \mathcal{L}(S / G) : \notag \\
                       &\delta(q_0, t) = q \wedge P(t) = P(s)\}.
\end{flalign}

% \textcolor{red}{BH: I suggest to use $\mathcal{E}_S(\alpha)$ directly for $\alpha\in P(\mathcal{L}(G))$}
%
% Thanks a lot

\begin{remark}\label{rem:uc_ctrl_pattern}
    % In DFAs, in accordance with the Ramadge-Wonham framework of \textup{\cite{ramadge1987supervisory}}, uncontrollable events are always enabled, regardless of the control patterns. 
    % That is, transitions labeled by uncontrollable events are guaranteed to be accessible regardless of the control policy.
    % %
    % In articles such as \cite{yin2016synthesis} and \cite{yin2018synthesis}, it is required that $\Sigma_{uc} \subseteq \gamma$, $\forall \gamma \in \Gamma$.
    % %
    % To highlight the controllability aspect of the system, we omit uncontrollable events from the control specifications $\gamma$ and $\Gamma$.
    % %
    % Their accessibility is treated as implicitly enabled in the remainder of this article.
    
    In the standard Ramadge-Wonham framework \cite{ramadge1987supervisory}, control patterns are often defined as subsets of $\Sigma$, and the uncontrollable events are always enabled by the supervisor regardless of the specific control pattern, i.e., $\Sigma_{uc}\subseteq \gamma$ for all $\gamma\in\Gamma$, as in \cite{yin2016synthesis,yin2018synthesis}.
    In this article, we adopt an equivalent convention by defining $\Gamma=2^{\Sigma_c}$, so that each $\gamma\in\Gamma$ records only the enabled controllable events.
    Uncontrollable events are treated as implicitly enabled throughout; equivalently, the actual enabled event set associated with $\gamma$ is $\gamma\cup\Sigma_{uc}$.
\end{remark}



% Let language $K \subseteq \mathcal{L}_m(G)$ be a non-prefix-closed language that represents the desired behavior to be achieved under control. $K$ is said to be:
% \begin{enumerate}[(i)]
%     \item controllable (w.r.t. DFA $G$ and uncontrollable event set $\Sigma_{uc}$) if $\overline{K} \Sigma_{uc} \cap \mathcal{L}(G) \subseteq \overline{K}$;
%     \item observable (w.r.t. DFA $G$, controllable event set $\Sigma_c$ and observable event set $\Sigma_o$) if for all $s \in \overline{K}$ and $\sigma \in \Sigma_c$,
%     such that $s \, \sigma \in \overline{K}$, $P^{-1}(P(s)) \, \sigma \cap \mathcal{L}(G) \subseteq \overline{K}$, where $P^{-1}(s) := \{ t \in \Sigma^* | P(t) = s \}$;
%     \item $\mathcal{L}_m(G)$-closed if $\overline{K} \cap \mathcal{L}_m(G) = K$. 
% \end{enumerate}
% %
% The following theorem shows that these three conditions together provide the necessary and sufficient conditions for the existence of a supervisor $S$ that exactly achieves the given language $K$.
% %
% \begin{theorem}[Controllability and Observability Theorem \cite{lin1988observability}]
%     Consider DES $G$, with controllable event set $\Sigma_c$ and observable event set $\Sigma_o$. Consider also $K \subseteq \mathcal{L}_m(G)$. 
%     There exists a supervisor $S : \mathcal{L}(G) \rightarrow \Gamma$ such that $\mathcal{L}_m(S /G) = K$ and $L(S / G) = \overline{K}$ if and only if
% \begin{enumerate}[(i)]
%     \item $K$ is controllable w.r.t. $\mathcal{L}(G)$ and $\Sigma_{uc}$;
%     \item $K$ is observable w.r.t. $\mathcal{L}(G)$, $\Sigma_o$ and $\Sigma_c$;
%     \item $K$ is $\mathcal{L}_m(G)$-closed.
% \end{enumerate}
% \end{theorem}

% If a given specification language cannot be exactly achieved,
% then the synthesis problem asks whether we can find a controllable, observable, and $\mathcal{L}_m(G)$-closed sublanguage of the specification language, and one that is ``as large as possible.''
% We would like to synthesize the non-blocking and maximally permissive, that is, the closed-loop system $S / G$ satisfies:
% \begin{enumerate}[(i)]
%     \item $\mathcal{L}_m(S /G) \subseteq K$ and $\mathcal{L}(S /G) = \mathcal{L}_m(S /G)$;
%     \item For any other a supervisor $S'$ satisfying (i), we have $\mathcal{L}_m(S / G) \not \subset \mathcal{L}_m(S' / G)$.
% \end{enumerate}
%
% Note that (ii) is usually relaxed in the scenario under partial observations, since observability is not preserved under union.


\section{Problem formulation}\label{sec:prob}

In this section, we introduce the notion of cognitive non-blocking based on the preliminaries and formulate the synthesis problems for DESs under this notion.
%
% TODO examples? 一个还是两个?
Then, we provide examples to illustrate the challenges of this problem.

\subsection{Cognitive non-blocking synthesis problem under partial observations}

In the previous section, Definition~\ref{def:nb} introduces the non-blocking property in DESs.
Under this property, every state that is \emph{accessible} from the initial state $q_0$ is also \emph{co-accessible} to the marked state $Q_m$.
%
The non-blocking supervisory control problem aims to synthesize a supervisor that selects admissible control decisions based solely on partial observations, such that every reachable sequence in the closed-loop system can be extended to reach the marked states. 
In this setting, the supervisor can restrict the behavior of the system by dynamically enabling or disabling the controllable events, thereby rejecting the reachability to the blocking states.

Formally, the non-blocking supervisory control problem under partial observation asks to find a supervisor $S: \mathcal{L}(G) \to \Gamma$ such that
\begin{flalign}\label{eq:nb_sc}
    \forall s \in \mathcal{L}(S / G), \exists t \in \Sigma^* : st \in \mathcal{L}_m(S / G),
\end{flalign}
or, in the closed-loop system $S / G$, all states accessible from the initial state are co-accessible to the marked states.
However, the non-blockingness only illustrates the capability of the states within DFA $G$ for arriving at a marked state.
In this sense, when the system is only partially observed, and the supervisor does not know the exact current state at each step, reaching a marked state in the actual system may still be insufficient for certifying, based solely on observations, that a marked state has indeed been reached.

Therefore, we would like to discuss the notion named \emph{cognitive non-blocking}, in which the task completion can be certified from partial observations, 
\begin{equation}
    \forall s \in \mathcal{L}(S/G), \exists t \in \Sigma^* : \mathcal{E}_S(st) \subseteq Q_m,
\end{equation}
%
% \textcolor{red}{This one is not the same as the one you used later, maybe explain why you use $\hat{\mathcal{E}}$, or directly introduce $\hat{\mathcal{E}}$ in your prelimi...}
%
% Thank you for your comment, we have added the explanation. We have refined the state estimate ... to ...
%
where $\mathcal{E}_S(st)$ is defined in \eqref{eq:esti:1}.
Intuitively, a DFA satisfies cognitive non-blocking if it can eventually certify, based on its observations, that it is in a marked state. 
That is, for every current behavior $s$, there exists a continuation $t$ such that the resulting state estimate $\mathcal{E}_S(st)$ consists entirely of marked states.
% It is straightforward to verify that cognitive non-blocking is a stronger property than classical non-blocking.

Before giving the problem formulation, we would like to discuss the physical interpretation behind this definition.
Cognitive non-blocking emphasizes the ability to \emph{certify} task completion based on observations.
%
However, if a marked state is reachable from a unmarked state via unobservable events, the system may enter the marked set silently, making it impossible for an external observer to determine whether the task has indeed been completed.
Therefore, interpreting reachability of marked states as task completion is meaningful only when such completion can be certified from observations.
% 
In many physical scenarios, such as a robot arriving at a goal region or a machine finishing an operation, completion is expected to be accompanied by some observable event.
%
To this end, we refine the state estimate $\mathcal{E}_S(\cdot)$ to $\hat{\mathcal{E}}_S(\cdot)$, introduce the notion of cognitive non-blocking, and then formally formulate the problem as follows.
%
\begin{definition}%[Cognitive Non-blocking]
    \label{def:cnb}
    Given DFA $G = (Q, \Sigma, \delta, q_0, Q_m)$ and a set of observable events $\Sigma_o \subseteq \Sigma$, we say DFA $G$ is \textit{cognitive non-blocking} with respect to observable event set $\Sigma_o$ if
    \begin{equation}\label{eq:cnb_1}
    \forall \alpha \in \mathcal{L}(S/G), \exists \omega \in \Sigma^* : \hat{\mathcal{E}}_S(\alpha \omega) \subseteq Q_m,
    \end{equation}
    % \textcolor{red}{BH: Also here, write as $\forall \alpha\in P(\mathcal{L}(G)), \exists \omega\in \Sigma_o^*$, so that this concept will look more like ``cognitive" since you define it directly on the observations.}
    %
    % thanks for your comment
    %
    where
    \begin{equation}\label{eq:est_o}
    \begin{aligned}
    \hat{\mathcal{E}}_S(\alpha) := \{ 
    &q \in Q \mid \exists \omega \in (\Sigma^*\Sigma_o) \cap \mathcal{L}(S / G) : \\ 
    &\delta(q_0, \omega) = q \wedge P(\omega) = P(\alpha) \}
    \end{aligned}
    \end{equation}
    %
    % TODO
    % 所有alpha要过一遍
    is the \emph{cognitive state estimate} associated with observation sequence $\alpha \in P(\mathcal{L}(G))$, i.e., the set of states that can be reached immediately after the immediate observation of $s$.
\end{definition}
%
\begin{problem}[Cognitive non-blocking synthesis problem]\label{prob:cnb}
    Given DFA $G = (Q, \Sigma, \delta, q_0, Q_m)$, a set of observable events $\Sigma_o \subseteq \Sigma$, and a set of controllable events $\Sigma_c \subseteq \Sigma$, synthesize the partial-observation supervisor such that
    \begin{enumerate}[(i)]
        \item The system is \textit{cognitive non-blocking} with respect to the marked states accessible via the observable events
        \begin{equation}
        \forall \alpha \in P(\mathcal{L}(S/G)),\
        \exists \omega \in \Sigma_o^* :        \hat{\mathcal{E}}_S(\alpha\omega)\subseteq Q_m.
        \end{equation}
        \item The supervisor is maximally permissive in the sense that for any other supervisor $S'$, we have
        $$
            \mathcal{L}(S/G) \not \subset \mathcal{L}(S'/G).
        $$
    \end{enumerate}
\end{problem}

\subsection{Challenge of synthesis problem}

% The synthesis of cognitive non-blocking is challenging because 

Non-blockingness is a state-based progress requirement in which every reachable state in the closed-loop system should have a continuation to the marked set $Q_m$.
Under partial observation, however, the supervisor does not know the exact state and must make control decisions based only on the observations generated so far.
Thus, even for the classical non-blocking objective, synthesis already has to account for indistinguishable states and uniform control decisions.
%
Cognitive non-blockingness imposes a stronger requirement.
It is not sufficient that the actual state can eventually reach $Q_m$.
Instead, after every generated behavior, there must exist a continuation whose observation allows the supervisor to certify that the current state is marked.
The progress target is not the marked set $Q_m$ itself, but a set of possible current states that is entirely contained in $Q_m$.
To this end, we provide Example~\ref{ex:nb:1} to illustrate why the traditional method cannot resolve the cognitive non-blocking synthesis problem.

\begin{example}\label{ex:nb:1}
Let us consider DFA $G$ in Fig. \ref{fig:dfa_grid_env} (b) with observable event set $\Sigma_o = \{o_1, o_2, o_3\}$, controllable event set $\Sigma_c = \{ c_1, c_2, o_3 \}$, and the marked state set
$Q_m = \{ q_3, q_5, q_6, q_{18} \}$.
This DFA is constructed from the motivating example of Fig.~\ref{fig:dfa_grid_env} (a). 
Backward reachability search (BRS) is a standard fixed-point computation used in supervisor synthesis \cite{thuijsman2020transformational, thuijsman2024reducing}.
For the satisfaction of classical non-blockingness, BRS starts from the marked set and repeatedly adds states that can reach the current retained set.
%
We show how to find the non-blocking supervisor $S$ using BRS.
Let the initial state set be
$
    V_0 = Q_m =\{ q_3, q_5, q_6, q_{18} \}.
$
% For a DFA under supervisory control, the controller can enable or disable controllable events to prevent blocking states from being accessible.
First, based on the transition relations
$$
\begin{aligned}
    q_1\xrightarrow{o_3}q_3    &\quad q_1\xrightarrow{o_2}q_5 \quad
    q_2\xrightarrow{o_2}q_6    \\
    q_8\xrightarrow{o_1}q_{18} &\quad q_{17}\xrightarrow{o_1}q_3,
\end{aligned}
$$
the backward step gives
$
    V_1 = V_0 \cup \{ q_1, q_2, q_8, q_{17} \}.
$
Continuing the same computation, we obtain
$$
\begin{aligned}
    V_2 &= V_1 \cup \{q_0,q_4,q_7\}   \\
    V_3 &= V_2 \cup \{q_{16},q_{19}\} \\
    V_4 &= V_3 \cup \{q_{12},q_{15}\} \\
    V_5 &= V_4 \cup \{q_{11}\}.
\end{aligned}
$$
Hence, the co-accessible fixed point is
$
    V^* = \bigcup_{i = 1}^5 P_5.
$
The states outside this fixed point are
$
    Q \setminus V^* = \{ q_9, q_{10}, q_{13}, q_{14} \}.
$
Note that $q_{13}$ and $q_{14}$ are unmarked deadlock states, and
$
    q_9\xrightarrow{o_2}q_{13}, \;
    q_{10}\xrightarrow{o_2}q_{14},
$
where $o_2$ is uncontrollable. 
Therefore, $q_9$ and $q_{10}$ must also be excluded.
%
Since all entries to $q_9$ and $q_{10}$ are controllable, the classical BRS
synthesizes a non-blocking supervisor by disabling
$$
\begin{aligned}
    q_7\xrightarrow{c_2}q_9     &\quad q_8\xrightarrow{c_2}q_{10} \\
    q_{11}\xrightarrow{c_1}q_9  &\quad q_{12}\xrightarrow{c_1}q_{10},
\end{aligned}
$$
such that the controlled system $S / G$ is non-blocking in the classical state-based sense.
%
However, following the control-observation sequence 
$
    \{o_3\}\,o_3\,(\{c_1\}\,o_1\,\{ \emptyset \}\,o_1)^\omega,
$
or
$
    \{o_3\}\,o_3\,\{c_1\}\,o_1\,(\{ \emptyset \}\,o_1\,\{c_1\}\,o_1\,)^\omega, \in (\Gamma \times \Sigma_o)^*,
$
one can easily check that the cognitive state estimate
$$
\begin{aligned}
    \hat{\mathcal{E}}_S(o_3) 
    &= \hat{\mathcal{E}}_S(o_3  \, (o_1 \, o_1))
     = \hat{\mathcal{E}}_S(o_3  \, (o_1 \, o_1)^2) \\
    &= \cdots = \{ q_3, q_4 \} \\
    %
    \hat{\mathcal{E}}_S(o_3 \, o_1) 
    &= \hat{\mathcal{E}}_S(o_3 \, (o_1 \, o_1)   \, o_1)
     = \hat{\mathcal{E}}_S(o_3 \, (o_1 \, o_1)^2 \, o_1)\\
    &= \cdots = \{ q_{17}, q_{18} \}.    
\end{aligned}
$$
With $q_3,q_{18}\in Q_m$ and $q_4,q_{17}\notin Q_m$, the supervisor cannot certify that the system has reached a marked state from these observations.
\end{example}


Example~\ref{ex:nb:1} further shows that the non-blocking supervisor synthesized by BRS may fail to be cognitively non-blocking because BRS reasons over  individual states rather than observation-compatible state estimates.
%
We also provide Example~\ref{ex:nb:2} to illustrate that this traditional method may fail for the synthesis of the cognitive non-blocking problem.

\begin{example}\label{ex:nb:2}
We now show that the non-blocking closed-loop system $S/G$ obtained in Example~\ref{ex:nb:1} is not necessarily cognitively non-blocking. 
The key point is that BRS checks the co-accessibility of individual states, whereas cognitive non-blockingness requires certification based on observations.
%
From the viewpoint of an observer, both $q_3$ and $q_4$ are non-blocking since $q_3$ is marked, and $q_4$ can reach the marked state $q_{18}$ through
$
    q_4 \xrightarrow{c_1}q_8 \xrightarrow{o_1}q_{18}.
$
%
Therefore, the classical BRS has no reason to remove either $q_3$ or $q_4$, and the observable controllable event $o_3$ may remain enabled.
%
However, after observing $o_3$, the supervisor cannot distinguish the two runs
$$
    q_0\xrightarrow{u_1}q_1\xrightarrow{o_3}q_3,
    \qquad
    q_0\xrightarrow{u_2}q_2\xrightarrow{o_3}q_4 .
$$
The set of possible current states is
$
    \{q_3, q_4\},
$
which does not satisfy the notion of cognitive non-blocking for $q_4 \not \in Q_m$.
%
Moreover, after the same successor observation $o_1$, the two compatible runs may evolve as
$$
    q_3 \xrightarrow{c_1}q_7 \xrightarrow{o_1}q_{17},
    \qquad
    q_4 \xrightarrow{c_1}q_8 \xrightarrow{o_1}q_{18}.
$$
Thus, the supervisor can only infer that the current state belongs to
$
    \{q_{17},q_{18}\},
$
which does not satisfy the notion of cognitive non-blocking for $q_{17} \not \in Q_m$.
%
Therefore, along the same observable behavior, there exist two indistinguishable DFA runs: one visits
$
    q_3 \longrightarrow q_{17},
$
and the other visits
$
    q_4 \longrightarrow q_{18}.
$
Although each individual state satisfies the classical non-blocking requirement, the observations mix marked and unmarked possibilities, and the subsequent observation still cannot resolve this ambiguity.
Consequently, along with this observable behavior, the resulting state estimate does not consist solely of marked states. 
Thus, task completion cannot be certified from observations, and the closed-loop system is not cognitively non-blocking.
\end{example}


Examples~\ref{ex:nb:1} and~\ref{ex:nb:2} together show that cognitive non-blocking synthesis is not merely a problem of retaining as many classically non-blocking states as possible. 
The BRS computation in Example~\ref{ex:nb:1} removes the classical blocking states and yields a non-blocking closed-loop system, while Example~\ref{ex:nb:2} shows that the same evolution may still contain observations that mix marked and unmarked possibilities. 
Consequently, a closed-loop system synthesized by the standard BRS may be non-blocking in the classical sense but fail to be cognitively non-blocking. 
This is why cognitive non-blocking synthesis must make control decisions based on the evolution of observation-compatible state sets, rather than applying the state-based synthesis methods directly.



\section{Bipartite cognitive controller}\label{sec:bcc}

In this section, we attempt to solve the cognitive non-blocking problem using information-state-based structures.
%
We first give the notion of information state and the corresponding update rules to represent both unobservable reach under control patterns and observable reach with respect to specific observations.
%
We then put forward the bipartite cognitive controller (BCC) that captures the microscopic evolution of the information states.


\subsection{Information state and update rules}

% TODO
% 1 不要用粗体，读者很难受
% 2 不要套旧有的壳
\begin{figure}[ht!]
    \centering
    \includegraphics[width=\linewidth]{Figs/BTS_only.pdf}
    \caption{Evolution of information states of DFA $G$ in Fig.~\ref{fig:dfa_grid_env} (b), where the orange rectangles are for Y-states, and the blue rounded rectangles are for Z-states, respectively, and the control pattern $\emptyset$ in Z-states means no events are enabled.}
    \label{fig:bts_only}
\end{figure}

Since we are dealing with partially observed systems, we define the notion of an information state as $X = 2^Q$.
We then define the Y-states $Y \subseteq X$ for possible reachable states under partial observations,
and Z-states $Z \subseteq X \times \Gamma$ for states of unobservable reach under supervisory control.
%
The \emph{unobservable reach} of the information states $x \in X$ under
the control pattern $\gamma \subseteq \Sigma$ is given by
\begin{flalign}
    UR_\gamma(x) := \{ &q' \in Q \mid \exists q \in x, \exists \omega \in ((\Sigma_{uc} \cup \gamma) \cap \Sigma_{uo})^* \cup \notag \\
    & \{ \varepsilon \} : q' = \delta(q, \omega) \}.
    \label{eq:ur}
\end{flalign}
The \emph{observable reach} of information states $x \in X$ under observable event $\sigma_o \in \Sigma_o$ is given by
\begin{equation}
    NX_{\sigma_o}(x):=\{q'\in Q\mid \exists q\in x : \delta(q,\sigma_o)=q'\}.
    \label{eq:nx}
\end{equation}
%
For simplicity, we use $\bm{q}$  to denote $Y \cup Z$.
We also denote the possible reachable states and the control pattern as $I(z) = x$ and $\Gamma(z) = \gamma$ for Z-state $z = (x, \gamma) \in Z$, respectively.

We provide the following example to illustrate how the information state evolves based on the rules of \eqref{eq:ur} and \eqref{eq:nx}, and the overall process for the evolution of information states in Fig.~\ref{fig:bts_only}.
%
\begin{example}\label{ex:is_update}
    Let us consider again DFA $G$ in Fig.~\ref{fig:dfa_grid_env} (b) with observable event set $\Sigma_o = \{o_1, o_2, o_3\}$ and controllable event set $\Sigma_c = \{ c_1, c_2, o_3 \}$. 
    %
    We demonstrate the evolution of the information states as follows.
    %
    For the initial state $y_0 = \{ q_0 \}$, since $u_1, u_2$ are uncontrollable and unobservable, the only available control pattern is $\gamma_0 = \{ o_3 \}$.
    Next, we have $x_1 = UR_{\gamma_0}(y_0) = \{ q_0, q_1, q_2 \}$ and $z_1 = (x_1, \gamma_1) = (\{ q_0, q_1, q_2 \}, \{ o_3 \})$, where states $q_0$, $q_1$, $q_2$ are accessible from state $q_0$ through the unobservable transitions $\varepsilon$, $u_1$, and $u_2$, respectively.
    %
    For information state $z_1$, the observable reach based on observation $o_3$ and $o_2$ is
    $y_1 = NX_{o_3}(x_1) = \{ q_3, q_4 \}$ and $y_2 = NX_{o_2}(x_1) = \{ q_5, q_6 \}$.
    Similarly, based on the possible unobservable transition $q_3 \xrightarrow{c_1} q_7 \xrightarrow{c_2} q_9$ and $q_4 \xrightarrow{c_1} q_8 \xrightarrow{c_2} q_{10}$, we can enable no event, namely $\emptyset$, or enable $\{ c_1 \}$, $\{ c_2 \}$, or $\{ c_1,c_2 \}$ at this step.
    Then, we have
    $x_4 = UR_{\gamma_2}(y_1) = \{ q_3, q_4, q_7, q_8 \}$ and
    $x_7 = UR_{\gamma_3}(y_1) = \{ q_3, q_4, q_7, q_8, q_9, q_{10} \}$,
    where $\gamma_2 = \{ c_1 \}$ and $\gamma_3 = \{ c_1, c_2 \}$.
    For Z-states $z_7 = (x_7, \gamma_3)$, we have
    $y_5 = NX_{o_2}(x_7) = \{ q_{13}, q_{14} \}$ and
    $y_3 = NX_{o_1}(x_7) = \{ q_{17}, q_{18} \}$.
\end{example}

\subsection{Construction of bipartite cognitive controller}

Based on the update rules of information states, we need a structure that can not only track the evolution of the controlled system based on the partial observations, but also track the microscopic structure within unobservable reaches and observable reaches.
%
The bipartite cognitive controller (BCC) structure fulfills this role by expanding information states and recording admissible control sequences.
Now, we give the formal definition of BCC as follows.
%
\begin{definition}\label{def:bcc}
Given a DFA $G$, its corresponding BCC is a tuple
$
\mathcal{C} = (\bm{\mathcal{Q}}, \Sigma_\mathcal{C}, \delta_\mathcal{C}, q^Y_0, \bm{\mathcal{Q}}_m)
$
where
\begin{enumerate}[(i)]
    \item $\bm{\mathcal{Q}} \subseteq \bm{\mathcal{Q}}^Y \dot{\cup} \bm{\mathcal{Q}}^Z$ is the set of microstates defined by
    $$
    \begin{aligned}
        \bm{\mathcal{Q}}^Y &= \{ (y, q) \in Y \times Q \mid y \in Y \wedge q \in y \} \\
        \bm{\mathcal{Q}}^Z &= \{ (z, q) \in Z \times Q \mid z \in Z \wedge q \in I(z) \},        
    \end{aligned}
    $$

    \item $\Sigma_\mathcal{C} = \Sigma \cup \Gamma$ is the set of transition labels.

    \item $\delta_\mathcal{C}: \bm{\mathcal{Q}} \times \Sigma_\mathcal{C} \rightarrow \bm{\mathcal{Q}}$ is the partial transition function, such that
    \begin{itemize}
        \item $\delta_\mathcal{C}((y,q_1),\gamma) = (z,q_2)$ if $q_1 = q_2$ and 
        $$
        I(z) = UR_{\gamma}(y) \wedge \Gamma(z) = \gamma,
        $$
        %
        \item $\delta_\mathcal{C}((z,q_1), \sigma_{uo}) = (z,q_2)$ if $\delta(q_1,\sigma_{uo})=q_2$ and $\sigma_{uo} \in (\Gamma(z)\cup\Sigma_{uc}) \cap \Sigma_{uo}$,
        %
        \item $\delta_\mathcal{C}((z,q_1),\sigma_o) = (y,q_2)$ if $\delta(q_1, \sigma_o)=q_2$, $\sigma_o \in (\Gamma(z) \cup \Sigma_{uc}) \cap \Sigma_o$, and
        $$
        y = NX_{\sigma_o}(I(z)).
        $$
    \end{itemize}

    \item $q^Y_0 = (\{q_0\}, q_0)$ is the initial state.

    \item $\bm{\mathcal{Q}}_m = \{ (y,q) \in \bm{\mathcal{Q}}^Y \mid y\subseteq Q_m, q\in y \}$ is the set of cognitively marked microstates.
\end{enumerate}
\end{definition}
%
In BCC $\mathcal{C}$, the first component of each state is its information state, whereas the second component indicates its DFA state. 
We use $q^Y = (y, q) \in \bm{\mathcal{Q}}_y$ and $q^Z = (z, q) \in \bm{\mathcal{Q}}_z$ to denote the individual microstates.
Given Y-state $y \in Y$ and Z-state $z \in Z$, its corresponding microstates, $\bm{\mathcal{Q}}_y$ and $\bm{\mathcal{Q}}_z$ are defined as 
$$
\begin{aligned}
    \bm{\mathcal Q}_y &= \{ (y,q) \in \bm{\mathcal Q}^Y \mid q \in y \} \\
    \bm{\mathcal Q}_z &= \{ (z,q) \in \bm{\mathcal Q}^Z \mid q \in I(z) \},
\end{aligned}
$$
%
respectively.
A \emph{run} of BCC $\mathcal{C}$ is a sequence of microstates
\begin{equation}\label{eq:bcc_run}
    \rho_\mathcal{C} = q_0^Y q_{1,1}^Z \cdots q_{1,m_1}^Z q_1^Y q_{2,1}^Z \cdots,
\end{equation}
and we denote the corresponding \emph{sequence} of control patterns and events for $\mathcal{C}$ as
$$
\bm{s} = tr_\mathcal{C}(\rho_\mathcal{C}) = \gamma_0 \sigma_{uo,1,1} \cdots \sigma_{uo,1,m_1} \sigma_{o,1} \gamma_1 \cdots \in \Sigma_\mathcal{C}^*.
$$
Here, the boldface Latin letters, such as $\bm{s}$, are used to explicitly record these sequences, distinguishing them from event sequences of DFAs.
%
Following the definition of the transition function $\delta$ for DFA $G$, we define the partial transition function
$$
(\bm{q}', q') = \delta_\mathcal{C}((\bm{q}, q), tr_\mathcal{C}(\rho_\mathcal{C})),
$$ 
such that $(\bm{q}, q)$ and $(\bm{q}', q')$ are the originating and ending microstates of the run $\rho_\mathcal{C}$, respectively.
We also define the full set of the runs of BCC as $\mathcal{R}_\mathcal{C}$.
%
Intuitively, the BCC structure is introduced to describe the ``game'' between the supervisor and the system under partial observation. 
%
It represents not only the control-observation evolution at the level of information states, but also the microscopic transitions within each information state.
The Y-states represent sets of observation-equivalent states and serve as decision-making instances for the supervisor, whereas the Z-states represent the resulting information states together with the selected decision $\Gamma(z)$, from which the system resolves the next observable event.
%
The transition from $q^Y \xrightarrow{\gamma} q^Z$ captures the unobservable reach under decision $\gamma$.
The transition $q^Z \xrightarrow{\sigma_{uo}} q^{Z\prime}$ captures the transitions within an unobservable reach, where $q^Z$ and $q^{Z\prime}$ share identical Z-state.
The transition $q^Z \xrightarrow{\sigma_o} q^{Y\prime}$ captures the observable reach after an event $\sigma_o \in \Sigma_o$, where $q^{Y\prime}$ consists of the states reachable via $\sigma_o$.
%
We use operator $R_{X} : \mathcal{R}_{\mathcal{C}} \rightarrow X^*$ to obtain the information state sequence from the BCC run $\rho_{\mathcal{C}}$, such that
\begin{flalign}
    R_{X}(\rho_\mathcal{C}) =  
    &y_0 z_0 y_1 z_1 \cdots \in X^* : \notag\\
    & \rho_\mathcal{C} = q_0^Y q_{1,1}^Z \cdots q_{1,m_1}^Z q_1^Y q_{2,1}^Z \cdots \wedge \ \notag\\
    &q_i^Y = (y_i, q_i), \ q_{i,j}^Z = (z_i, q'_{i,j}), \ \forall i,j \in \mathbb{N},
\end{flalign}
%
We use $\iota = R_{X}(\rho_{\mathcal{C}})$ to denote the resulting sequence of information state.
%
We also define operator $R_{X}^{-1} : X^* \rightarrow 2^{\mathcal{R}_{\mathcal{C}}}$  for information state sequence $\iota = R_X(\rho_\mathcal{C})= y_0 z_0 y_1 z_1 \cdots$, such that
$$
\begin{aligned}
    R_{X}^{-1}(\iota) := \{ &\rho_{\mathcal{C}}' \in \mathcal{R}_\mathcal{C} \mid 
        \rho_{\mathcal{C}}' = (y_0, q_0) (z_0, q_{1,1}') \cdots (z_0, q_{1,m}') \\  
            &(y_1, q_{y,1}') \cdots : q_{1,1}',\ldots,q_{1,m}' \in I(z_0), \ \\
            &q_{y,1}' \in y_1, \ldots \}.
\end{aligned}
$$
Note that the sequence of information states $\iota$ may correspond to multiple BCC runs in $R_{X}^{-1}(\rho_{\mathcal{C}})$.
We define the available control pattern at Y-state $y$ as
\begin{equation}\label{eq:ctrl_y}
    H(y) := \{ \gamma \in \Gamma \mid \exists q \in y, \delta_{\mathcal{C}}((y, q), \gamma)! \}.
\end{equation}
%
% 下到上
% \highlight{HERE WE NEED TO REFINE THE STATE SEQUENCES RATHER THAN RUNS, JUST SEE $\{ q_{15}, q_{16} \}$ can reach $\{q_0\}$ via $o_3$ and $o_1$.}
%
The language generated by BCC $\mathcal{C}$ is defined by
\begin{equation}
    \mathcal{L}(\mathcal{C}) := \{ tr_{\mathcal{C}}(\rho_{\mathcal{C}}) \in \Sigma_{\mathcal{C}}^* \mid \delta_{\mathcal{C}}(q^Y_0, tr_{\mathcal{C}}(\rho_{\mathcal{C}}))! \}.
    \label{eq:lang_bcc}
\end{equation}
Since $\mathcal{C}$ contains both events and control patterns, we define the projection of events
$
    P_\Sigma: \Sigma_{\mathcal{C}}^* \to \Sigma^*
$
for BCC $\mathcal{C}$ recursively by
\begin{equation}\label{eq:seq_proj_bcc}
P_\Sigma(\alpha \eta)
=
\begin{cases}
    P_\Sigma(\alpha) \eta, & \eta \in \Sigma,\\
    P_\Sigma(\alpha),      & \eta \in \Gamma,
\end{cases}
\end{equation}
with $P_\Sigma(\varepsilon)=\varepsilon$,
$\forall \alpha \in \Sigma_{\mathcal{C}}^*$ and $\eta \in \Sigma_{\mathcal{C}}$.
Following the projection $P_{\Sigma}$ in \eqref{eq:seq_proj_bcc}, the projection onto control patterns is defined recursively as
\begin{equation}\label{eq:proj_bcc_ctrl}
P_\Gamma(\alpha \eta) =
\begin{cases}
P_\Gamma(\alpha) \eta, & \eta \in \Gamma,\\
P_\Gamma(\alpha), & \eta \in \Sigma,
\end{cases}
\end{equation}
%
with $P_\Gamma(\varepsilon)=\varepsilon$.
Finally, define the projected language of the event sequences induced by $\mathcal{C}$ as
\begin{equation}
    \mathcal{L}_{\Sigma}(\mathcal{C})
    :=
    P_\Sigma(\mathcal{L}(\mathcal{C}))
    \subseteq \Sigma^* .
    \label{eq:proj_lang_bcc}
\end{equation}
Since $\mathcal{L}_{\Sigma}(\mathcal{C})$ is generated by an automaton-based structure, it is naturally prefix-closed.


\begin{algorithm}[ht!]
\caption{Construction of BCC}
\label{alg:bcc}
\KwIn{DFA $G=(Q,\Sigma,\delta,q_0,Q_m)$ with event partitions
$\Sigma=\Sigma_c\dot{\cup}\Sigma_{uc}=\Sigma_o \dot{\cup} \Sigma_{uo}$\;}
\KwOut{BCC $\mathcal{C} = (\bm{\mathcal{Q}},\Sigma_{\mathcal{C}},
\delta_{\mathcal{C}},q^Y_0,\bm{\mathcal{Q}}_m)$\;}

$\Gamma \gets 2^{\Sigma_c}$\;
$y_0 \gets \{q_0\}$, $q^Y_0 \gets (y_0,q_0)$, $\Sigma_{\mathcal{C}} \gets \Sigma \cup \Gamma$\;
$\bm{\mathcal{Q}}\gets \{(y_0,q_0)\}$, $\delta_{\mathcal{C}}\gets \emptyset$\;
$Y\_stack\gets \{y_0\}$, $Y\_visited\gets \{y_0\}$\;

\While{$Y\_stack\neq \emptyset$}{
    Pop $y$ from $Y\_stack$\;

    % \tcp{Add all Y-microstates associated with $y$}
    \ForEach{$q\in y$}{
        Add $(y,q)$ to $\bm{\mathcal{Q}}$\;
    }

    % \tcp{Enumerate all control patterns for all-inclusive construction}
    \ForEach{$\gamma\in\Gamma$}{
        $x\gets UR_{\gamma}(y)$, $z\gets (x,\gamma)$\;

        \ForEach{$q\in x$}{
            %
            % \tcp{Add all Z-microstates}
            Add $(z,q)$ to $\bm{\mathcal{Q}}$\;
            %
            % \tcp{Y-Z decision transitions}
            \If{$q \in y \cap x$}{
                Add transition $((y,q), \gamma, (z,q))$ to $\delta_{\mathcal{C}}$\;
            }
            %
            % \tcp{Internal unobservable transitions inside the same Z-state}
            \ForEach{$\sigma_{uo} \in (\gamma \cup \Sigma_{uc})\cap\Sigma_{uo}$ $\wedge$ $\delta(q,\sigma)!$}{
                    $q'\gets \delta(q, \sigma_{uo})$\;
                    Add $(z,q')$ to $\bm{\mathcal{Q}}$\;
                    Add transition $((z,q), \sigma_{uo}, (z,q'))$ to $\delta_{\mathcal{C}}$\;
            }
        }

        % \tcp{Observable exits from the Z-state to new Y-states}
        \ForEach{$\sigma_o\in(\gamma \cup \Sigma_{uc})\cap\Sigma_o$}{
            $y'\gets NX_{\sigma_o}(x)$\;
            \If{$y'\neq\emptyset$}{

                \ForEach{$q'\in y'$}{
                    Add $(y',q')$ to $\bm{\mathcal{Q}}$\;
                }

                \ForEach{$q\in x$ $\wedge$ $\delta(q,\sigma_o)!$}{
                        $q'\gets \delta(q,\sigma_o)$\;
                        Add transition $((z,q),\sigma_o,(y',q'))$ to $\delta_{\mathcal{C}}$\;
                }

                \If{$y'\notin Y\_visited$}{
                    Push $y'$ to $Y\_stack$\;
                    $Y\_visited\gets Y\_visited \cup \{y'\}$\;
                }
            }
        }
    }
}

$\bm{\mathcal{Q}}_m \gets
\{(y,q)\in\bm{\mathcal{Q}} \mid y\subseteq Q_m,\ q\in y\}$\;

\Return $\mathcal{C}$\;
\end{algorithm}
%
We give Algorithm~\ref{alg:bcc} for the construction of a BCC from the original DFA $G$.
% 
Algorithm~\ref{alg:bcc} follows a breadth-first search (BFS) structure.
It iteratively finds all the successor Z-states from the given Y-state based on all available control pattern, and finds the successor Y-states from all observable reaches of the new Z-states.
Meanwhile, whenever a new information state is generated, the algorithm also constructs the associated microstates and transition relations induced by the underlying DFA transitions.
%
In Lines 5--10, the algorithm explores Y-states in a BFS manner. For each Y-state $y$, it adds all corresponding Y-microstates $(y,q)$ with $q\in y$ to the BCC, and enumerates control patterns to compute successor Z-states through the operator $UR_\gamma(\cdot)$. 
The corresponding Y-Z micro-transitions, which correspond to the selected control decision, are then added to the BCC in Line 14.
%
In Lines 11--18, the algorithm expands the newly generated Z-state by adding all Z-microstates $(z,q)$ with $q\in I(z)$ and recording the internal unobservable transitions within the same Z-state.
%
In Lines 19--29, the algorithm computes successor Y-states using the operator $NX_{\sigma_o}(\cdot)$ for observable events with nonempty reachable sets, and adds the corresponding Y-microstates and Z-Y transitions to the BCC.
%
We use Example~\ref{ex:bcc} to illustrate the construction of BCC.
%
\begin{example}\label{ex:bcc}
    Consider again the DFA $G$ in Fig.~\ref{fig:dfa_grid_env} (b). 
    We construct the corresponding BCC using Algorithm~\ref{alg:bcc}, with the resulting structure shown in Fig.~\ref{fig:bcc_overall}, where orange rectangles and blue rounded rectangles denote Y-states and Z-states, respectively. 
    For clarity, we simplify the notation $(\bm{q}, q)$ to $q$ and use different colors for the transition relations to highlight distinct evolutions, thereby illustrating the structural properties of each microstate.
\end{example}

% 
% 然后需要说明为什么不足够解决，参考一下博涵的聊天记录
% 原文中的NB是直接用AIC做的
% Z状态是deadlock-free是for all states的(公式(7))，而Y-状态是live的
% 这里似乎是不是Z的条件是一致的，似乎我们只要加强Y状态即可？
% 那么我的问题变成了是不是只用AIC就可以解决这个问题 (就是动机这关怎么过)
%   答案是:
%   Z是不变的，但是就算判断live也要用ICS (BCC)
%   如果没有 ICS/BCC 结构，你就无法直接判断 Y-state 是否 live，因为普通 NB-AIC 的 live 判断是针对真实状态，而 cognitive live 需要追踪信息状态 evolution（即观察到的 projection 下的可能状态集合）。
%
% 对于某个 Y-state y（对应 NB-AIC 中的信息状态集合），它是 cognitive live，当且仅当对于每个 x∈y 的可能真实状态，都存在一条控制决策序列，使得经过一系列允许的事件后，observer 能确认系统进入了 fully marked Y-state。
%
% 3. ICS/BCC 结构的必要性
% ICS/BCC（Information Consistency Structure / Bipartite Control Consistency）结构本质上是：
% 展开信息状态：对每个 Y-state，追踪在 observable projection 下，可能的 DFA states 集合 evolution。
% 可控决策序列：记录哪些控制决策保证从该信息状态出发可以到达 fully marked Y-state。
% 递归判定 live：通过 BFS/DFS 或 fixed-point，判断每个信息状态是否有至少一条控制序列使 cognitive estimate 最终落入 fully marked Y-state。
% 因此：
% 没有 ICS/BCC，就无法考虑信息状态 evolution，只能用 DFA state 判定 live；
% cognitive live 必须依赖 ICS/BCC 或等价结构，否则 supervisor 无法保证从 observer 角度看非阻塞。
% 换句话说：
%       cognitive live \neq = plant live (原 NB-AIC live)
% 只有通过信息状态展开（ICS/BCC），才能把 live 条件提升到 observer 可验证的层面。

%
% 这个模板caption中\bm{q}显示不出来
\begin{figure*}[t!]
    \centering
    \includegraphics[width=0.615\linewidth]{Figs/bcc_overall.pdf}
    \caption{Bipartite cognitive controller constructed from the DFA $G$ in Fig.~\ref{fig:dfa_grid_env} (b). We simplify the notations $(y,q)$ and $(z,q)$ to $q$ to better illustrate the structural properties of each microstate. Orange rectangles represent Y-states, while blue rounded rectangles represent Z-states. Transitions in different colors are used to distinguish different evolutions. Moreover, the information states marked by red crosses are infeasible Z-states, cognitively blocking Y-states, or inconsistent Y- or Z-states, which should be removed to achieve cognitive non-blocking. 
    The source code is available at \url{https://github.com/trigger1996/des_common}.}   % TODO, 最好拆成几个库
    \label{fig:bcc_overall}
\end{figure*}


\subsection{Properties of bipartite cognitive controller}

The BCC represents the system evolution at two levels: the information-state level and the DFA-state level. 
At the information-state level, it captures the evolution induced by partial observations and control decisions. 
At the DFA-state level, it records the detailed transitions among individual DFA states, including internal transitions within unobservable reaches and observable transitions between the microstates of Y- or Z-states.


% TODO
% $((\bm{q}, q), \eta, (\bm{q}', q'))$, 这个label符号要统一
Consider a BCC $\mathcal{C}$, where each microstate is of the form $(\bm{q}, q)$ with $\bm{q} \in X$ and $q \in Q$, and each transition is of the form $((\bm{q}, q), \eta, (\bm{q}', q'))$. 
Define
\begin{flalign}\label{eq:pre_post}
    Pre((\bm{q}, q))  = \{ &(\bm{q}', q') \in \bm{\mathcal{Q}} \mid \exists \eta \in \Sigma_{\mathcal{C}}, \nonumber \\ 
    &(\bm{q}, q) \in \delta_{\mathcal{C}}((\bm{q}', q'), \eta) \} \nonumber \\
    %
    Post((\bm{q}, q)) = \{ &(\bm{q}', q') \in \bm{\mathcal{Q}} \mid \exists \eta \in \Sigma_{\mathcal{C}}, \nonumber \\ 
    &(\bm{q}', q') = \delta_{\mathcal{C}}((\bm{q}, q), \eta) \}
\end{flalign}
as the sets of predecessors and successors of $(\bm{q}, q)$, respectively. 
For a set of states $\bm{\mathcal{Q}}_{s} \subseteq \bm{\mathcal{Q}}$, we use
\begin{flalign}\label{eq:pre_post:2}
    Pre(\bm{\mathcal{Q}}_{s})  := &\bigcup_{(\bm{q}, q) \in \bm{\mathcal{Q}}_{s}} Pre((\bm{q}, q)) \nonumber \\
    Post(\bm{\mathcal{Q}}_{s}) := &\bigcup_{(\bm{q}, q) \in \bm{\mathcal{Q}}_{s}} Post((\bm{q}, q))
\end{flalign}
to denote the sets of predecessors and successors, respectively. 
%
For $(\bm{q}, q), (\bm{q}', q') \in \bm{\mathcal{Q}}$ and $\bm{\mathcal{Q}}_s \subseteq \bm{\mathcal{Q}}$, we write $(\bm{q}', q') \leadsto_{\bm{\mathcal{Q}}_s} (\bm{q}, q)$ if there exists a nonempty path from $(\bm{q}', q')$ to $(\bm{q}, q)$,
$$
(\bm{q}', q') (\bm{q}_{1}, q_1) (\bm{q}_{2}, q_2) \cdots (\bm{q}_{n}, q_n) (\bm{q}, q),
$$
such that $(\bm{q}_{1}, q_1), \ldots, (\bm{q}_{n}, q_n) \in \bm{\mathcal{Q}}_s$.
A state $(\bm{q}, q) \in \bm{\mathcal{Q}}$ is called a \emph{deadlock} if 
\begin{equation}\label{eq:bcc_dl}
    Post((\bm{q}, q)) = \emptyset.
\end{equation}
A nonempty set $\bm{\mathcal{Q}}_l \subseteq \bm{\mathcal{Q}}$ is called a \emph{livelock} if $\bm{\mathcal{Q}}_l$ is a nontrivial strongly connected component, and $Post(\bm{\mathcal{Q}}_l) \subseteq \bm{\mathcal{Q}}_l$, i.e.,
\begin{flalign}\label{eq:bcc_ll}
    &\bm{\mathcal{Q}}_l = \emptyset \wedge
    \forall (\bm{q},q), (\bm{q}',q') \in \bm{\mathcal{Q}}_l : 
    ((\bm{q},q) \leadsto_{\bm{\mathcal{Q}}_l} \nonumber \\
    & (\bm{q}',q') \wedge
      (\bm{q}',q') \leadsto_{\bm{\mathcal{Q}}_l} (\bm{q},q)) \wedge
    Post(\bm{\mathcal{Q}}_l) \subseteq \bm{\mathcal{Q}}_l.   
\end{flalign}

Based on the above notations, we analyze the properties of the BCC, illustrating the consistency of its evolution with the original DFA, the completeness of the control representation, and its construction complexity in terms of both time and space.
%
We first clarify a structural restriction on Y-states. 
Recall that Y-states represent observation-equivalent states under supervisory control, and thus transitions between Y-microstates must preserve the observation structure. In particular, for any $y \in Y$ and any $(y, q), (y, q')$ with $q, q' \in y$, we have
$
((y, q), \eta, (y, q')) \notin \delta_{\mathcal{C}}, \forall \eta \in \Sigma_{\mathcal{C}}.
$
%
Building on this restriction, we describe the internal structure of Z-states $z$ and their corresponding Z-microstates $(z, q) \in \bm{\mathcal{Q}}_z$.
A Z-microstate is called an \emph{entry Z-microstate} if all its predecessors $Pre((z,q))$ belong to $\bm{\mathcal{Q}}^Y$, and if $Post((z,q)) \cap \bm{\mathcal{Q}}^Y \neq \emptyset$, it is called an \emph{exit Z-microstate}. 
Note that a Z-microstate can be both an entry and an exit microstate.
%
For an entry Z-microstate $(z,q)$ and any predecessor Y-microstate $(y,q)$, the identity transition in \eqref{eq:ur} implies $q = \delta(q,\varepsilon)$, and thus $q \in I(z)\cap y$. 
This shows that Z-microstates inherit their information state from Y-microstates under unobservable transitions.
%
Moreover, transitions within a Z-state are induced by unobservable reachability $UR_\gamma(\cdot)$ starting from its predecessor Y-states.
%
Therefore, reachability within a Z-state is consistent with this inheritance structure. 
Specifically, for any entry Z-microstate $(z,q)$ and its predecessor Y-microstate $(y, q)$, every $q' \in I(z)$ is reachable via a string 
$
s \in ((\Gamma(z)\cup\Sigma_{uc})\cap\Sigma_{uo})^*
$ 
such that
$$
\delta_{\mathcal{C}}((y,q_y),\Gamma(z))=(z,q_y),
\quad
\delta_{\mathcal{C}}((z,q_y), s)=(z, q').
$$
%
We give the following proposition to illustrate the consistency between the DFA $G$ and the BCC $\mathcal{C}$.
%
\begin{proposition}\label{prop:bcc_dfa_consistency}
Given a run 
$$
\begin{aligned}
\rho_{\mathcal{C}} = &(y_0, q_0) (z_1, q_{1,1}) \cdots (z_1, q_{1,m_1}) (y_1, q_1) \cdots \\
                     &(y_i, q_i) (z_{i+1}, q_{i+1,1}) \cdots (z_{i+1}, q_{i+1,m_{i+1}}) (y_{i+1}, q_{i+1}) \cdots
\end{aligned}
$$
in the BCC $\mathcal{C}$, there exists a corresponding sequence
% 这个s是对的, 因为这是 DFA 序列
$
s = \xi_0 \sigma_{o,0} \xi_1 \sigma_{o,1} \cdots \xi_i \sigma_{o,i} \cdots
$
in the closed-loop system $S/G$ such that:
\begin{enumerate}[(i)]
    \item Each $\xi_i \in ((\Gamma(z_i) \cup \Sigma_{uc}) \cap \Sigma_{uo})^*$ is an unobservable subsequence corresponding to the unobservable reach $(z_{i}, q_{i,1}) \cdots (z_{i}, q_{i,m_{i}}) (y_{i}, q_{i})$;
    \item For each subsequence $(z_i, q_{i,1}) \cdots (z_i, q_{i,m_i})$, there exist events $\sigma_{i,j} \in \Sigma_{uo}$ such that
    $
    (q_{i,j}, \sigma_{i,j}, q_{i,j+1}) \in \delta;
    $
    \item For each transition $((z_i, q_{i,m_i}), \sigma_{o,i}, (y_{i+1}, q_{i+1}))$, we have
    $
    (q_{i,m_i}, \sigma_{o,i}, q_{i+1}) \in \delta.
    $
\end{enumerate}
\end{proposition}
%
\begin{proof}
The result follows directly from the construction of $\delta_{\mathcal{C}}$.
%
First, every Z-Z transition satisfying $(z, q_{i,j+1}) = \delta_{\mathcal{C}}((z, q_{i,j}), \sigma)$, $i,j \in \mathbb{N}$ is defined only if $\sigma \in (\Gamma(z) \cup \Sigma_{uc}) \cap \Sigma_{uo}$ and $(q_{i,j}, \sigma, q_{i,j+1}) \in \delta$. 
Thus, each subsequence $(z_i, q_{i,1}) \cdots (z_i, q_{i,m_i})$ corresponds to a sequence $\xi_i \in ((\Gamma(z_i) \cup \Sigma_{uc}) \cap \Sigma_{uo})^*$ satisfying (ii).
%
Second, every Z-Y transition $(y_{i+1}, q_{i+1}) = \delta_{\mathcal{C}}((z_{i}, q_{i,j}), \sigma_o)$ is defined only if $\sigma_o \in \Sigma_o$ and $(q, \sigma_o, q') \in \delta$ of DFA $G$, which establishes (iii).
%
Finally, every Y-Z transition $((y_i, q_i), \gamma_{i+1}, (z_{i+1}, q_{i+1,1}))$ is defined by $\delta_\mathcal{C}$ such that $\gamma_{i+1} = \Gamma(z_{i+1})$.
This determines the set of admissible events $(\Gamma(z_{i+1}) \cup \Sigma_{uc})$ that govern the subsequent unobservable evolution within $z_{i+1}$, and thus constrains the corresponding subsequence $\xi_{i+1}$.
Concatenating all such unobservable subsequences $\xi_i$ and observable events $\sigma_{o,i}$ yields a valid execution $s$ in the closed-loop system $S / G$, which proves (i).
\end{proof}
%
As a direct consequence of Proposition~\ref{prop:bcc_dfa_consistency}, 
for any $\alpha \in \mathcal{L}(\mathcal{C})$, if
$
\delta_{\mathcal{C}}(q^Y_0,\alpha)=(\bm{q},q),
$
then
$
    \delta(q_0,P_\Sigma(\alpha)) = q.
$
In particular, we have
$$
    \mathcal{L}_{\Sigma}(\mathcal{C})
    =
    P_\Sigma(\mathcal{L}(\mathcal{C}))
    \subseteq
    \mathcal{L}(G).
$$
That is, every DFA sequence projected from the BCC is feasible and belongs to the language $\mathcal{L}(G)$.
%
Moreover, we provide the time complexity of Algorithm~\ref{alg:bcc} and the space complexity of BCC as follows.


\begin{proposition}
Given a DFA $G$ with $|Q|$ states and $|\Sigma|$ events, the complexity of Algorithm~\ref{alg:bcc} in time will be dominated by
$$
O(2^{|Q| +|\Sigma|} \cdot (|Q|\cdot |\Sigma|))
$$
\end{proposition}


\begin{proof}
Algorithm~\ref{alg:bcc} performs a breadth-first search over generated Y-states. Since each Y-state is a subset of $Q$, the number of repetitions of Line 5 is at most
$
    |Y| \leq 2^{|Q|}.
$
%
In Line 7, the Y-microstates will be added to BCC, and the repetition will be at most  $O(|Q|)$.
%
In Line 9, for each generated Y-state $y$, the algorithm enumerates all control patterns
$
    \gamma\in\Gamma=2^{\Sigma_c},
$
such that there are at most
$
    |\Gamma| = 2^{|\Sigma_c|}
$
iterations of the inner loop.
%
We next estimate the cost for a fixed pair $(y, \gamma)$ in Line 10, which corresponds to the admissible control pattern at state $y$.
Denote the set of admissible observable and unobservable reachable events as
$
    \mathcal{U}_\gamma = (\gamma \cup \Sigma_{uc})\cap\Sigma_{uo},
$
and
$
    \mathcal{O}_\gamma = (\gamma \cup \Sigma_{uc})\cap\Sigma_o .
$
First, the algorithm computes
$
    x = UR_\gamma(y).
$
Using a standard graph-search implementation over the DFA transition table, this requires visiting at most $|Q|$ states and checking at most 
$| \mathcal{U}_\gamma | \leq |\Sigma|$ events at each state. Thus, the computation complexity of $UR_\gamma(y)$ is
$
    O(|Q|\cdot |\Sigma|).
$
%
After $x = I(z)$ is obtained, the algorithm adds the corresponding Z-microstates and records the internal unobservable transitions inside the same Z-state. 
This amounts to scanning the subgraph induced by $x$ under events in $\mathcal{U}_\gamma$.
Since $|x|\leq |Q|$ and the number of such internal transitions is at most
$|Q|\cdot |\Sigma|$, this step costs
$
    O(|Q|\cdot |\Sigma|).
$
%
Finally, in Line 20, for each $\sigma_o\in \mathcal{O}_\gamma$, it computes
$
    y' = NX_{\sigma_o}(x)
$
and records the corresponding Z-Y micro-transitions. 
Since $|\mathcal{O}_\gamma|\leq |\Sigma|$ and each computation scans at most $|Q|$ states, this step also costs
$
    O(|Q|\cdot |\Sigma|).
$
The total cost of Algorithm~\ref{alg:bcc} is
$$
O(|Y| \cdot 2^{|\Sigma|} \cdot (|Q| +3(|Q|\cdot |\Sigma|))),
$$
which will be dominated by
$
O(2^{|Q| +|\Sigma|} \cdot (|Q|\cdot |\Sigma|)).
$
\end{proof}



\begin{proposition}\label{prop:state_bcc}
Consider a DFA $G$ with $|Q|$ states and $|\Sigma|$ events. For the corresponding BCC $\mathcal{C}$ constructed by Algorithm~\ref{alg:bcc}, the number of microstates and transitions satisfy
$$
\begin{aligned}
    |\bm{\mathcal{Q}}|
    &\leq
    |Q|(2^{|Q|}+2^{|Q|+|\Sigma|}) \\
    |\delta_{\mathcal{C}}|
    &\leq
    |Q|\,2^{|Q|+|\Sigma|}(1+|\Sigma|),
\end{aligned}
$$
respectively.
\end{proposition}

\begin{proof}
Each BCC microstate is of the form $(\bm{q},q)$, where $\bm{q}$ is either a Y-state or a Z-state, and $q\in Q$ is a DFA state contained in the corresponding information state.
Thus, the maximal number of states satisfies
$
    |\bm{\mathcal{Q}}|
    \leq
    (|Y|+|Z|)\cdot |Q|.
$
Since each Y-state is a subset of $Q$, we have
$
    |Y|\leq 2^{|Q|}.
$
As for the Z-states in the form of
$
    z=(x,\gamma),
$
where $x\subseteq Q$ and $\gamma\in 2^{\Sigma_c}$. 
Therefore, the maximal number of Z-states is
$
    |Z|\leq 2^{|Q|}\cdot 2^{|\Sigma_c|}
    =
    2^{|Q|+|\Sigma_c|},
$
and it follows that
$$
    |\bm{\mathcal{Q}}|
    \leq
    |Q|(2^{|Q|}+2^{|Q|+|\Sigma_c|}).
$$

Next, we calculate the maximal number of transitions. 
The BCC contains three types of micro-transitions: Y-Z transitions, internal Z-Z transitions, and Z-Y transitions.
%
For Y-Z transitions, for each Y-state $y$, the control pattern satisfies $\gamma \in 2^{\Sigma_c}$.
Thus, the number of Y-Z transitions via control pattern $(y,q)\xrightarrow{\gamma}(z,q)$ is bounded by
$$
    |Y|\cdot 2^{|\Sigma_c|}\cdot |Q|
    \leq
    |Q|\,2^{|Q|+|\Sigma_c|}.
$$
%
Given a Z-state $z=(x,\gamma)$, each microstate $(z,q)$ has at most one successor
under each event $\sigma\in\Sigma$, since the DFA $G$ is deterministic.
Therefore, the total number of internal Z-Z transitions and Z-Y transitions over all Z-states is bounded by
$
    |Z|\cdot |Q|\cdot|\Sigma|.
$
Since
$
    |Z|\leq 2^{|Q|+|\Sigma_c|},
$
we obtain
$
    |Z|\cdot |Q|\cdot|\Sigma|
    \leq
    |Q|\cdot|\Sigma|\cdot 2^{|Q|+|\Sigma_c|}.
$
%
Combining the above bounds, we obtain
$$
\begin{aligned}
    |\delta_{\mathcal{C}}|
    &\leq
    |Q|\,2^{|Q|+|\Sigma_c|}
    +
    |Q|\cdot|\Sigma|\cdot 2^{|Q|+|\Sigma_c|} \\
    &=
    |Q|\,2^{|Q|+|\Sigma_c|}(1+|\Sigma|).
\end{aligned}
$$

Since $|\Sigma_c|\leq|\Sigma|$, we further have
$$
    |\bm{\mathcal{Q}}|
    \leq
    |Q|(2^{|Q|}+2^{|Q|+|\Sigma|}),
$$
and
$$
    |\delta_{\mathcal{C}}|
    \leq
    |Q|\,2^{|Q|+|\Sigma|}(1+|\Sigma|).
$$
\end{proof}






\section{Solution to Cognitive non-blocking synthesis problem}\label{sec:solu}
% 
% 我后面需要什么，我前面需要什么
% 
%
% 1 consistency of BCC 作为下一章节的第一个定理
% 2 what is potentially feasible Z-state?
% 3 what is cognitive live Y-state
% 4 how to synthesis CNB
% 5 Algorithm 2
% 6 correctness of Algorithm 2 (有必要用单源单汇路径的结论吗？)
% 7 

In this section, we leverage the definition of consistency for BCCs.
By pruning those inconsistent states in BCC, the sublanguage generated by the BCC can be guaranteed to be cognitively non-blocking.
%

\subsection{Synthesis algorithm}\label{subsec:supervisor_induction}

To synthesize cognitive non-blocking sublanguages, we first identify the
target microstates characterized by \eqref{eq:cnb_1}.
%
Recall that in Definition~\ref{def:bcc}, the cognitively marked states $\bm{\mathcal{Q}}_m$ are defined as the set of Y-microstates such that for each $q^Y \in \bm{\mathcal{Q}}_m$, its corresponding Y-state $y_m \in Y$ consists entirely of marked states. 
For each such cognitively marked Y-state $y_m$, we denote its corresponding Y-microstate set as $\bm{\mathcal{Q}}_{y,m}$.
%
Intuitively, according to \eqref{eq:cnb_1}, a target Y-state corresponds to an information state in which all observation-consistent states are marked. 
Therefore, regardless of which actual DFA state has been reached under the partial observation, it can be certified that the system has reached a marked state.
%
Meanwhile, we define the \emph{consistency} of a BCC $\mathcal{C}$ as follows.
%
\begin{definition}\label{def:con_bcc}
Given a BCC $\mathcal{C}$, we say it is \emph{consistent} if
\begin{enumerate}[(i)]
    \item For each Z-microstate $(z, q)$, all possible observable transitions from its information state $z \in Z$ are defined, 
    $
    \forall (z, q) \in \bm{\mathcal{Q}}^Z, \forall \sigma_o \in \Gamma(z)\cap\Sigma_o:\;
    (\delta(q,\sigma_o)!) \Rightarrow \delta_{\mathcal{C}}((z, q), \sigma_o)!;    
    $
    \item For each Y-microstate $(y, q)$, at least one control decision is defined for its information state $y \in Y$, 
    $
    \forall (y, q) \in \bm{\mathcal{Q}}^Y, y \in Y : H(y) \neq \emptyset.
    $
\end{enumerate}
\end{definition}
%
Note that in Definition~\ref{def:bcc} (iii), only Z-states can be admissible successors for Y-states in BCC $\mathcal{C}$.
%
Next, we give the following two definitions for properties of Y-states and Z-states based on the reachability of their microstates, which will be further utilized for the notion of cognitive non-blocking.
%
% \begin{definition}\label{def:z_pf}
%     We say a Z-state $z \in Z$ is \emph{potentially feasible} if, for each of its microstates $q^Z \in \bm{\mathcal{Q}}_z$, it does not lead to a deadlock \eqref{eq:bcc_dl} or a livelock \eqref{eq:bcc_ll} consisting solely of the Z-microstates $\bm{\mathcal{Q}}_z$ of $z$. 
% \end{definition}
\begin{definition}\label{def:z_pf}
    We say a Z-state $z \in Z$ is \emph{potentially feasible} if, for each of its microstates $q^Z \in \bm{\mathcal{Q}}_z$, there does not exist a deadlock \eqref{eq:bcc_dl} reachable from $q^Z$, nor does there exist a livelock \eqref{eq:bcc_ll} contained entirely in $\bm{\mathcal{Q}}_z$.
\end{definition}
%
% 这里有一点很关键，这个Y-state必须能到同一个mark，不能(y, q)到一个marked, (y, q')到另一个mark
\begin{definition}\label{def:y_cl}
    We say a Y-state $y$ is \emph{cognitively live} with respect to the set of marked Y-microstates $\bm{\mathcal{Q}}_m$ if and only if there exists a cognitively marked Y-state $y_m \subseteq Q_{m}$ and a information-state sequence
    $
    \iota = y z_0 y_1 z_1 \cdots z_k y_m,
    $
    such that all Y-microstates $(y, q)$ satisfy
    \begin{flalign}\label{eq:bcc_y_cl}
        &\forall q \in y, \exists q_m \in y_m : \rho_{\mathcal{C}} \in R_X^{-1}(\iota) : \notag \\
        &(y_m, q_m) = \delta_{\mathcal{C}}((y,q), tr_{\mathcal{C}}(\rho_{\mathcal{C}})),
    \end{flalign}
    where $z_k$, $y_m$, $k \in \mathbb{N}, m \in \mathbb{N}_{\geq 0}$ are potentially feasible Z-microstates and other cognitively live Y-states, respectively.
\end{definition}
%
Intuitively, a Y-state is cognitively live if all its microstates can reach the microstates of cognitively a marked Y-state via the runs $R_X^{-1}(\iota)$ corresponding to the information state sequence $\iota$ of the potentially feasible Z-microstates and other cognitively live Y-states.

The microstate-level of Definition~\ref{def:z_pf} and \ref{def:y_cl} are defined upon their corresponding information states.
If a Z-state is potentially feasible, then all of its microstates are potentially feasible, and vice versa.
If a Y-state is said to be cognitively live, then all its Y-microstates are cognitively live.
%
% 第二个没有 vice versa
%
We give the following example to illustrate in detail the property of potential feasibility and cognitive liveness in a BCC.

\begin{figure}[ht!]
    \centering
    \includegraphics[width=0.975\linewidth]{Figs/bcc_routes_overall.pdf}
    \caption{Illustration of several evolutions in the BCC $\mathcal{C}$ shown in Fig.~\ref{fig:bcc_overall}. 
    In particular, subfigure~(b) depicts two possible branches associated with DFA state $q_{17}$. 
    For simplicity, the corresponding Z-state is represented by either $(\{ q_{15}, q_{16}, q_{17} \}, \{ c_1 \})$ or $(\{ q_{15}, q_{16}, q_{17} \}, \{ c_1, o_3 \})$. 
    The key difference between these two branches is whether event $o_3$ is enabled. 
    When $o_3$ is enabled, the Y-microstate $(\{q_0\}, q_0)$ becomes accessible from the Z-state of $q_{17}$; otherwise, it does not. This distinguishing transition is highlighted in orange.}
    \label{fig:bcc_routes}
\end{figure}
%
\begin{example}\label{ex:y_z_prop}
Let us consider again the DFA $G$ in Fig.~\ref{fig:dfa_grid_env} (b) and the corresponding BCC $\mathcal{C}$ in Fig.~\ref{fig:bcc_overall}.
%
For the evolution shown in Fig.~\ref{fig:bcc_routes} (a), the Z-state $z_6 = (\{ q_{13}, q_{14} \}, \emptyset)$ is not potentially feasible, since both of its Z-microstates $(z_6, q_{13})$ and $(z_6, q_{14})$ are deadlock states.
%
Meanwhile, the Y-state $y_5 = \{ q_{13}, q_{14} \}$ is not cognitively live, because its only available transition leads to the deadlock Z-state $z_6$.
The Y-states $y_1 = \{ q_{3}, q_{4} \}$ and $y_3 = \{ q_{17}, q_{18} \}$ are not cognitively live either.
Since the only available marked Y-state is $y_2 = \{ q_5, q_6 \}$, there is no available sequence from the Y-microstates  $(y_1, q_3)$, $(y_1, q_4)$, $(y_3, q_{17})$, and $(y_3, q_{18})$ to the marked microstates $(y_2, q_5)$ and $(y_2, q_6)$. 
Thus, \eqref{eq:bcc_y_cl} is not satisfied.
%
As for the evolution shown in Fig.~\ref{fig:bcc_routes} (b), the Y-state $y_4 = { q_{15}, q_{16} }$ is cognitively live.
For the state sequence $\iota = y_4 z_{11} y_0 z_0$, with $z_{11} = (\{ q_{15}, q_{16}, q_{17}\}, \{c_1, o_3\})$, $y_0 = {q_0}$, and $z_0 = (\{ q_0, q_1, q_2 \}, \emptyset)$, we have the corresponding BCC runs
$
\rho_{\mathcal{C}, 1} = (y_4, q_{15}) \xrightarrow{\{c_1, o_3\}} (z_{11}, q_{15}) \xrightarrow{c_1} (z_{11}, q_{17}) \xrightarrow{o_3} (y_0, q_0) \xrightarrow{\emptyset} (z_0, q_0) \xrightarrow{u_1} (z_0, q_1) \xrightarrow{o_2} (y_2, q_5)
$
and
$
\rho_{\mathcal{C}, 2} = (y_4, q_{16}) \xrightarrow{\{c_1, o_3\}} (z_{11}, q_{16}) \xrightarrow{o_1} (y_0, q_0) \xrightarrow{\emptyset} (z_0, q_0) \xrightarrow{u_2} (z_0, q_2) \xrightarrow{o_2} (y_2, q_6),
$
both of which end at microstates associated with the same cognitively marked Y-state $y_2$.
%
In this sense, one can also check that $(y_0, q_0)$ is cognitively live.
Moreover, the Z-states $z_{11} = (\{ q_{15}, q_{16}, q_{17}\}, \{c_1, o_3\})$ and $z_5 = (\{q_5, q_6, q_{11}, q_{12}\}, \{ c_1 \})$ are potentially feasible, since their exit Z-microstates, $(z_5, q_{11})$, $(z_5, q_{12})$, $(z_{11}, q_{16})$, and $(z_{11}, q_{17})$, have cognitively live successors.
However, if the controllable event $o_3$ is not enabled at $y_4$, then the run $\rho_{\mathcal{C}, 1}$ is no longer feasible.
Therefore, $y_4$ is no longer cognitively live.
%  
% That is, by observing the subsequence $(o_3 o_2)^\omega$ or $(o_1 o_2)^\omega$, the observer can certify that the system arrives at a marked state.
%
% 这个问题在于15、16到q_0的trace不一样，一个o_1，一个o_3, 所以不能是run到run，必须是state sequence到run
%
\end{example}


%
Example~\ref{ex:y_z_prop} shows that microstates must be retained or removed as whole units. 
Since an information state represents observation-equivalent DFA states, retaining only part of its microstates would break the observation-based consistency. 
% Thus, if any associated microstate violates Definition~\ref{def:con_bcc}, or prevents the corresponding Y- or Z-state from satisfying Definitions~\ref{def:z_pf} or~\ref{def:y_cl}, the entire information state $y$ or $z$ should be removed.

%
% 1   考虑Y和Z是耦合的，而Z直接影响系统的可行性，Y只需要保证有通路即可
% 2   首先标记那些不可行的Z-state，通过Z-microstates判断
% 2.2 此时infeasible Z可以分为两类：一类是all所有结束状态为deadlock
% 2.3                             一类是存在exit states
% 3   首先删除2.2和2.3
% 4   再根据consistency倒着删，此时一边确定系统是否consistent，一边确定Y是否cognitively live，直到系统CNB
%
%
% 我知道最后的结果就是系统会很破碎，但是：不行的状态就是不行
To synthesize a supervisor that guarantees that the closed-loop system is cognitive non-blocking, we need to prune those states that do not satisfy the properties of potential feasibility, cognitive liveness, or consistency.
%
Although the accessibility of Y-states and Z-states is coupled, which determines the non-blocking property of the system during its evolution, 
any Z-microstate that is a deadlock or livelock within an unobservable reach is clearly infeasible. 
Thus, we start by identifying Z-states that are not potentially feasible. 
Specifically, we first detect the Z-microstates corresponding to deadlocks and livelocks, as defined in \eqref{eq:bcc_dl} and \eqref{eq:bcc_ll}, and then prune the associated Z-states from the BCC.
%
The removal of infeasible Z-states may affect the consistency and cognitive liveness of predecessor Y-states. 
For each predecessor Y-state $y \in Y$ of a pruned Z-state, 
if $y$ is no longer cognitively live, i.e., \eqref{eq:bcc_y_cl} does not hold for some $q \in y$, or it violates Definition~\ref{def:con_bcc} (i) — then $y$ should also be pruned.
%
After updating Y-states, some remaining Z-states may become inconsistent. 
We recursively prune inconsistent Z-states and their predecessor non-cognitively-live Y-states until no further removals are necessary and the system becomes consistent. 

We summarize this procedure in Algorithm~\ref{alg:cnb_syn} and illustrate it with Example~\ref{ex:cnb_removal}.
%
% TODO
% 这部分的正确性的concern的担忧主要来自于删除顺序，因为这个问题比原有的IS-based问题在删除过程中会有更强的耦合性，所以得看怎么删
% 最好的办法是过一遍代码检查一下
Note that the resulting Y-states after the pruning procedure are not guaranteed to be accessible from the initial state. 
In Line 7, the operator $Accessible(\cdot)$ extracts the accessible component from the initial state $q^Y_0$.
%
% 这里只要有q_0就能把所有状态连起来, 所以不用union of component
%
It is required that, for each microstate in the returned BCC $\mathcal{C}_c$, the corresponding Y- or Z-state is \emph{consistent}, i.e., all microstates associated with a given information state are included, so that the structural integrity of the BCC is preserved.  
%
For the remainder of the BCC $\mathcal{C}_c$, if there exists at least one Y-microstate corresponding to a cognitively marked state in $\bm{\mathcal{Q}}_{m}$, then $\mathcal{C}_c$ represents a cognitive non-blocking supervisor with respect to the marked states.
Otherwise, no cognitive non-blocking supervisor exists, and the algorithm correctly returns ``No solution''.
%

\begin{algorithm}[t!]
\caption{Synthesis of Cognitive Non-Blocking BCC}
\label{alg:cnb_syn}
\KwIn{BCC $\mathcal{C} = (\bm{\mathcal{Q}}, \Sigma_\mathcal{C}, \delta_\mathcal{C}, q^Y_0, \bm{\mathcal{Q}}_m)$}
\KwOut{Cognitive non-blocking BCC $\mathcal{C}_c$}

\BlankLine

% \tcp{Step 1: Prune infeasible Z-states}
Identify all Z-microstates in $\bm{\mathcal{Q}}^Z_l \subseteq \bm{\mathcal{Q}}$ that are deadlock or livelock using \eqref{eq:bcc_dl} and \eqref{eq:bcc_ll}\;
Prune the Z-states $z$, $(z, q) \in \bm{\mathcal{Q}}^Z_l$ from $\bm{\mathcal{Q}}$\;

\Do{BCC $\mathcal{C}$ not consistent}{

    Prune the Z-states that violate Definition~\ref{def:con_bcc} (i)\;
    Prune the Y-states that are not cognitively live or violate Definition~\ref{def:con_bcc} (ii)\;

}
$\mathcal{C}_c = (\bm{\mathcal{Q}}_c, \Sigma_\mathcal{C}, \delta_{\mathcal{C}, c}, q^Y_0, \bm{\mathcal{Q}}_{m,c}) \gets Accessible(\mathcal{C})$\;
\If{$q_0^Y\notin \bm{\mathcal Q}_c$ or $\bm{\mathcal Q}_{m,c}=\emptyset$}{
    \Return No solution\;
}
\Return BCC $\mathcal{C}_c$\;
\end{algorithm}


\begin{example}\label{ex:cnb_removal}
    Let us consider the DFA in Fig.~\ref{fig:dfa_grid_env} (b) and the corresponding BCC in Fig.~\ref{fig:bcc_overall} again.
    We use Algorithm~\ref{alg:cnb_syn} to prune the states that do not satisfy the properties of potential feasibility and cognitive liveness, and then we make this system consistent.
    %
    First, following the result of Example~\ref{ex:y_z_prop}, one can easily check that the blocking Z-states $z_2$, $z_3$, and $z_6$ do not satisfy the definition of potential feasibility, and should be pruned.
    We also prune $z_{11}$ and $z_{12}$ because their microstates $(z_{11}, q_{19})$ and $(z_{12}, q_{15})$ are deadlock states.
    %
    Next, based on the definition of consistency, we iteratively remove $y_5$, $z_7$, and $z_8$.
    %
    As for Y-states $y_1$ and $y_3$, one can easily check that there is no feasible sequence that originates from the microstates of $y_1$ or $y_2$ and ends at cognitively marked state $(y_2, q_5)$ or $(y_3, q_6)$, which violates the requirement of cognitive liveness.
    We need to prune these states, and then we prune $z_4$ for consistency.
    Finally, based on Definition~\ref{def:con_bcc} (i), all enabled observations must be defined for Z-states, so the Z-state $z_1$ should be removed.
    %
    The removed states are marked by red crosses in Fig.~\ref{fig:bcc_overall}, and the remaining structure forms the resulting consistent cognitive non-blocking BCC.
\end{example}


In the remainder of this article, to simplify the expression, we use $\mathcal{C}$ instead of $\mathcal{C}_c$ to denote the consistent cognitive non-blocking BCC output by Algorithm~\ref{alg:cnb_syn}.

% TODO
% 因为这么删确实会把原系统弄得很破碎
\begin{remark}
The BCC obtained by Algorithm~\ref{alg:cnb_syn} may become ``fragmented'' after pruning. This is natural since cognitive non-blocking is stronger than classical non-blocking.
It requires not only reachability of marked states but also observation-based certification of task completion for those marked states.
Therefore, some non-blocking states may still be removed. 
% For supervisor synthesis, we simply restrict attention to the accessible component of the pruned BCC, namely, the part reachable from the initial state $q^Y_0$.
\end{remark}


%
% TODO
% 其实这图我感觉没啥用但是也就是好看
\begin{figure*}
    \centering
    \includegraphics[width=0.695\linewidth]{Figs/conceptual_proof.pdf}
    \caption{Conceptual illustration of the structural soundness in Theorem~\ref{theo:cnb}. 
    A retained BCC run projects to a feasible DFA run through $P_\Sigma$. For an arbitrary reachable microstate, a cognitively live Y-state admits a continuation to a marked Y-state $y_m\subseteq Q_m$, while a potentially feasible Z-state admits an observable exit to such a live Y-state. 
    Red-crossed components indicate cognitive-blocking or infeasible branches that are excluded from any BCC satisfying Definitions~\ref{def:con_bcc}, \ref{def:z_pf}, and \ref{def:y_cl}.}
    \label{fig:conceptual_proof}
\end{figure*}

\subsection{Synthesized supervisor and correctness}

In this subsection, we show how the cognitive non-blocking BCC $\mathcal{C}$ induces a supervisor and establishes the correctness of the synthesis algorithm.
%
Recall that the projection from the BCC language to the sequence of control patterns is
$
P_{\Gamma}:\Sigma_{\mathcal{C}}^* \rightarrow \Gamma^*,
$
and the set of event sequences generated by BCC is denoted as $\mathcal{L}_{\Sigma}(\mathcal{C})$.
For a BCC run
$
\rho_\mathcal{C}
=
(\bm{q}_0,q_0)(\bm{q}_1,q_1)\cdots(\bm{q}_n,q_n)\cdots,
$
the corresponding event sequence and control-pattern sequence are
$
    P_\Sigma(tr_\mathcal{C}(\rho_\mathcal{C}))\in\Sigma^*,
$
and
$
    P_\Gamma(tr_\mathcal{C}(\rho_\mathcal{C}))\in\Gamma^*,
$
respectively.
%
%
Note that a conventional partial-observation supervisor $S: P(\mathcal{L}(G)) \rightarrow \Gamma$ assigns a control pattern $\gamma \in \Gamma$ based on the given observation sequence. 
In the BCC framework, we use $S_{\mathcal{C}} : P(\mathcal{L}_{\Sigma}(\mathcal{C})) \rightarrow P_\Gamma(\mathcal{L}(\mathcal{C}))$ in place of $S$ as a history-based representation of the supervisor induced by the BCC. 
Specifically, $S_{\mathcal{C}}$ maps each projected event sequence to the sequences of admissible control patterns from the BCC $\mathcal{C}$, rather than the point-wise control pattern. 
For each observation sequence $\alpha \in P(\mathcal{L}_{\Sigma}(\mathcal{C}))$, with a slight abuse of notation, define
%
% \begin{flalign}
%     S_{\mathcal{C}}(\alpha) :=\{&\gamma\in\Gamma\mid \exists \bm{s} \in \mathcal L(\mathcal C): P(P_\Sigma(\bm{s}))=\alpha, \notag \\
%     &\bm{s}\gamma\in\mathcal L(\mathcal C)\}
% \end{flalign}
% 
\begin{flalign}\label{eq:bcc_sc:1}
%
% 试过了，这个就是一行放不下
S_{\mathcal{C}}(\alpha) = \{ &\nu \in \Gamma^* \mid\ \exists \bm{s} \in \mathcal{L}(\mathcal{C}) : \notag \\
&P(P_\Sigma(\bm{s})) = \alpha, \nu = P_\Gamma(\bm{s}) \}.
%
% //
% S_{\mathcal{C}}(\alpha) :=\{&\gamma\in\Gamma\mid \exists \bm{s} \in \mathcal{L}(\mathcal C): P(P_\Sigma(\bm{s}))=\alpha, \notag \\
% &\bm{s}\gamma\in\mathcal{L}(\mathcal C)\}.
\end{flalign}
%
%
That is, $S_{\mathcal{C}}(\alpha)$ collects all sequences of control patterns from those BCC runs whose projected event sequences can generate the same observation $\alpha$.
%
The closed-loop language induced by $S_{\mathcal{C}}$ is then defined by the event projections of the same BCC runs
\begin{flalign}\label{eq:bcc_sc:2}
\mathcal{L}(S_{\mathcal{C}}/G)
:=&
\{\, s\in\Sigma^* \mid
\exists \bm{s} \in\mathcal{L}(\mathcal C) :  s=P_\Sigma(\bm{s}),  \notag\\
& \; \; P_\Gamma(\bm{s}) \in S_{\mathcal{C}}(P(s))
\,\} \notag \\
=&
P_\Sigma(\mathcal{L}(\mathcal C))
=
\mathcal{L}_{\Sigma}(\mathcal C).
\end{flalign}
%
By \eqref{eq:bcc_sc:1}, the condition $P_\Gamma(\bm{s}) \in S_{\mathcal{C}}(P(s))$ means that the projection onto control pattern for the same BCC run $\bm{s}$ is admissible under the observation $P(s)$, while $s = P_\Sigma(\bm{s})$ extracts the corresponding event projection.
Thus, \eqref{eq:bcc_sc:1} gives the interpretation of the control sequence for a BCC run, whereas \eqref{eq:bcc_sc:2} extracts its event-sequence interpretation.
%
By Proposition~\ref{prop:bcc_dfa_consistency}, every such projected event sequence is a feasible execution of the DFA $G$.
Hence, we can apply $S_\mathcal{C}$ as $S$ such that $\mathcal{L}(S_{\mathcal{C}} / G)$ can be used as an equivalent closed-loop representation of the controlled DES, $\mathcal{L}(S/G)$.
%
Now, we give the following two theorems to show the structural characterization of cognitive non-blockingness. 
%
% Soundness of the algorithm
\begin{theorem}\label{theo:cnb}
Given a DFA $G = (Q, \Sigma, \delta, q_0, Q_m)$ and a corresponding BCC
$
{\mathcal{C}} = ({\bm{\mathcal Q}},\Sigma_{\mathcal{C}},
\delta_{\mathcal{C}}, q^Y_0, \bm{\mathcal{Q}}_m),
$
suppose that all microstates of ${\mathcal{C}}$ are accessible from $q^Y_0$ and satisfy Definitions~\ref{def:con_bcc}, \ref{def:z_pf}, and \ref{def:y_cl}.
The supervisor $S$ induced by ${\mathcal{C}}$ guarantees that the closed-loop system $S/G$ is cognitive non-blocking.
\end{theorem}
%
\begin{proof}
The intuition of the proof is illustrated in Fig.~\ref{fig:conceptual_proof}. The proof consists of two steps.
First, we use the BCC--DFA consistency to identify the closed-loop behavior $S/G$ with the event projection of the BCC language.
Second, we show that the retained BCC structure is co-accessible, in the cognitive sense, to a marked Y-state.

By the definition of the supervisor induced by $\mathcal{C}$, the closed-loop language generated by $S / G$ coincides with the event projection of the
BCC language, i.e.,
\begin{equation}\label{eq:proof_cnb_lang_equiv}
    \mathcal{L}(S/G) = \mathcal{L}_{\Sigma}(\mathcal C)
                     = P_{\Sigma}(\mathcal{L}(\mathcal C)).
\end{equation}
Moreover, by Proposition~\ref{prop:bcc_dfa_consistency}, every BCC run projects to a feasible execution of the DFA. Hence, for any $\bm{s} \in \mathcal{L}(\mathcal C)$ such that
$
    \delta_{\mathcal{C}}(q^Y_0, \bm{s}) = (\bm{q}, q),
$
we have
$
    \delta(q_0, P_\Sigma(\bm{s})) = q.
$
Therefore, it is sufficient to prove the cognitive non-blocking condition on the BCC and then project the result to the DFA.

Since $\mathcal{C}$ is accessible, for an arbitrary microstate $(\bm{q}, q) \in \bm{\mathcal{Q}}$, there exists a sequence $\bm{s}_{q_0, q}$ in $\mathcal{C}$ such that
$$
(\bm{q}, q) = \delta_{\mathcal{C}}((y_0, q_0), \bm{s}_{q_0, q}),
$$
%
and the corresponding event sequence is denoted by
$P_\Sigma(\bm{s}_{q_0, q}) = s$.
%
We next show that this reachable microstate admits a cognitively marked continuation. 
By the structural co-accessibility guaranteed in the BCC, there exists a sequence $\bm{t}_{q, q_m}$ such that
$$
(y_m, q_m) = \delta_{\mathcal{C}}((\bm{q}, q), \bm{t}_{q, q_m}),
$$
and the corresponding event sequence is denoted by $P_\Sigma(\bm{t}_{q, q_m}) = t$.
%
If $(\bm{q},q)$ is a Z-microstate, then Definition~\ref{def:z_pf} guarantees that it can reach a Y-microstate of a successor Y-state.
%
For a Y-microstate, \eqref{eq:bcc_y_cl} guarantees that, whenever its Y-state is cognitively live, all corresponding Y-microstates can reach a cognitively marked state through continuations following the same sequence of Y- and Z-states.
%
Moreover, the bipartite structure preserves observation consistency among microstates with the same first component.
Hence, for any microstate $(\bm{q}, q')$ that shares identical observations with $(\bm{q}, q)$, such that
$$
\begin{aligned}
    (\bm{q}, q') &= \delta_{\mathcal{C}}((y_0, q_0),   \bm{s}_{q_0, q'}) \\
    (y_m,   q_m) &= \delta_{\mathcal{C}}((\bm{q}, q'), \bm{t}_{q', q_m}),
\end{aligned}
$$
with
$
P(P_\Sigma(\bm{s}_{q_0, q'})) = P(P_\Sigma(\bm{s}_{q_0, q}))
$
and
$
P(P_\Sigma(\bm{t}_{q', q_m})) = P(P_\Sigma(\bm{t}_{q, q_m})).
$
Consequently, for any event sequence 
$s \in \mathcal{L}_{\Sigma}(\mathcal{C})$, there exists a continuation 
$t \in \Sigma^*$ such that the BCC-level cognitive state estimate induced by the observation $P(st)$ satisfies
%
$$
\begin{aligned}
%\widehat{\bm{\mathcal{E}}}_{\mathcal{C}}(P(st)) := 
\{ &(\bm{q}, q) \in \bm{\mathcal{Q}} \mid 
\exists \bm{r} \in \mathcal{L}(\mathcal{C}) :  
r = P_{\Sigma}(\bm{r}) \in \Sigma^*\Sigma_o, \\
& \delta_{\mathcal{C}}(q_0^Y,\bm{r}) = (\bm{q},q),\
P(r) = P(st) \}
\subseteq \bm{\mathcal{Q}}_m .
\end{aligned}
$$
%
%Here, $\bm r$ is a BCC sequence, while $r=P_\Sigma(\bm r)$ is its projected DFA event sequence. 
Hence, the above inclusion states that all BCC microstates consistent with the same observation $P(st)$ are cognitively marked.
%
By \eqref{eq:proof_cnb_lang_equiv} and Proposition~\ref{prop:bcc_dfa_consistency}, the above BCC-level inclusion can be projected back to the corresponding closed-loop DFA $S_{\mathcal{C}}/G$. Since every microstate in $\bm{\mathcal{Q}}_m$ is of the form $(y,q)$ with $y\subseteq Q_m$, its DFA component satisfies $q\in Q_m$. 
Therefore, for any observation sequence $\alpha \in P(\mathcal{L}(S_{\mathcal{C}}/G))$, there exists a continuation $\omega \in \Sigma_o^*$ such that the corresponding cognitive state estimate satisfies 
$
\hat{\mathcal{E}}_{S_{\mathcal{C}}}(\alpha\omega) \subseteq Q_m .
$
This completes the proof.
\end{proof}

% Soundness: 只要系统证明了某件事，那么这件事在语义上确实成立。
% Completeness: 只要某件事实际上成立，系统就一定能把它证明/表示出来。
\begin{theorem}\label{theo:cnb_complete}
Given a DFA $G = (Q, \Sigma, \delta, q_0, Q_m)$, suppose that the closed-loop system $S / G$ satisfies the notion of cognitive non-blocking.
Then, the corresponding BCC
$
{\mathcal{C}} = ({\bm{\mathcal Q}},\Sigma_{\mathcal{C}},
\delta_{\mathcal{C}}, q^Y_0, \bm{\mathcal{Q}}_m)
$
constructed for $S / G$ satisfies Definitions~\ref{def:con_bcc},
\ref{def:z_pf}, and~\ref{def:y_cl}.
\end{theorem}

\begin{proof}
By Definition~\ref{def:cnb}, for any $s \in \mathcal{L}(S / G)$ with $q = \delta(q_0, s)$, there exists a continuation $t$ such that $st \in (\Sigma^*\Sigma_o) \cap \mathcal{L}(S / G)$ and, for any $t' \in (\Sigma^*\Sigma_o) \cap \mathcal{L}(S / G)$,
\begin{equation}\label{eq:theo_comp:1}
    P(t') = P(st) \Rightarrow \delta(q_0, t') \in Q_m.
\end{equation}
%
We now examine how the reachable state $q$ appears in the corresponding information-state evolution. If $q=q_0$, then $q$ belongs to the initial Y-state $y_0$, and the following Y-state argument applies directly.

First, suppose that the incoming transition of $q$ is observable and can be written as $q' \xrightarrow[]{\sigma_o} q$. Then there exists an information state $z' = (x', \gamma)$ with $q' \in x'$ such that
$$
NX_{\sigma_o}(x') = y, \quad q \in y,
$$
i.e., $q$ belongs to a Y-state $y$.
%
By \eqref{eq:theo_comp:1}, after the continuation $t$, all closed-loop executions observation-equivalent to $st$ end in marked states. Since observation-equivalent executions are collected into the same successor Y-state, the Y-state $y_m'$ reached after $P(st)$ satisfies
$
y_m' \subseteq Q_m.
$
Thus, $y_m'$ is cognitively marked, and $y$ is cognitively live.
Hence, Definition~\ref{def:y_cl} holds.

Second, suppose that the incoming transition of $q$ is unobservable and can be written as $q' \xrightarrow[]{\sigma_{uo}} q$. Then there exists a Y-state $y'$ and a control pattern $\gamma$ such that
$$
x = UR_{\gamma}(y') \wedge q \in x,
$$
i.e., $q$ belongs to the Z-state $z = (x, \gamma)$ and to the Z-microstate $(z,q)$.
%
If $z$ contained a deadlock microstate, then the corresponding closed-loop sequence would admit no continuation to a later observation certifying \eqref{eq:theo_comp:1}. If $z$ contained a livelock consisting only of microstates of the same Z-state, then the system could stay within the unobservable reach of $z$ forever without reaching a successor Y-state, again contradicting \eqref{eq:theo_comp:1}. Therefore, no reachable Z-state in $\mathcal C$ contains such deadlock or livelock microstates, and Definition~\ref{def:z_pf} holds.

It remains to verify consistency. Since ${\mathcal{C}}$ is induced by the closed-loop behavior of a partial-observation supervisor, every reachable Y-state has at least one control decision selected by $S$. 
Moreover, for each reachable Z-state, every feasible observable event enabled under its control pattern must be represented by a successor Y-state; otherwise, some possible
closed-loop evolution under the same observation would be missing from $\mathcal C$. 
Hence, Definition~\ref{def:con_bcc} also holds.
%
Consequently, the BCC constructed for the cognitive non-blocking closed-loop system satisfies Definitions~\ref{def:con_bcc}, \ref{def:z_pf}, and~\ref{def:y_cl}.
\end{proof}

Intuitively, a closed-loop system is cognitively non-blocking if and only if its behavior can be represented by an accessible BCC whose behavior satisfies the requirements of consistency, potential feasibility, and cognitive liveness.
%
We also give the following corollary for the correctness of Algorithm~\ref{alg:cnb_syn}.
The algorithm returns a nonempty BCC exactly when the cognitive non-blocking synthesis problem is solvable. 
If a nonempty BCC remains after pruning, it directly induces a cognitive non-blocking supervisor; conversely, if such a supervisor exists, then its closed-loop behavior admits a BCC representation satisfying the same structural requirements, and hence the pruning procedure cannot eliminate all states.

\begin{corollary}\label{cor:alg_cnb_sound}
Algorithm~\ref{alg:cnb_syn} returns a nonempty BCC rather than ``No solution'' if and only if there exists a partial-observation supervisor $S$ such that the closed-loop system $S/G$ is cognitively non-blocking.
%
Moreover, if Algorithm~\ref{alg:cnb_syn} returns a nonempty BCC $\mathcal{C}$, then the supervisor $S_{\mathcal{C}}$ induced by $\mathcal{C}$ guarantees that $S_{\mathcal{C}}/G$ is cognitive non-blocking.
\end{corollary}

\begin{proof}
We first prove the ``only if'' part of this corollary.
Suppose that Algorithm~\ref{alg:cnb_syn} returns a nonempty BCC $\mathcal{C}$. 
By construction, the algorithm removes all Z-states that are not potentially feasible, all Y-states that are not cognitively live, and all Y- or Z-states that violate the consistency requirement. It then extracts the accessible component from the initial microstate $q^Y_0$. 
Hence, $\mathcal{C}$ is accessible from $q^Y_0$ and satisfies Definitions~\ref{def:con_bcc}, \ref{def:z_pf}, and~\ref{def:y_cl}. 
By Theorem~\ref{theo:cnb}, the supervisor $S_{\mathcal{C}}$ induced by $\mathcal{C}$ guarantees that $S_{\mathcal{C}}/G$ is cognitive non-blocking. 
Therefore, there exists a partial-observation supervisor enforcing cognitive non-blocking.

We next prove the ``if'' part. 
Suppose that there exists a partial-observation supervisor $S$ such that the closed-loop system $S/G$ is cognitively non-blocking. 
Let ${\mathcal{C}}'$ be the corresponding BCC constructed for $S / G$. 
By Theorem~\ref{theo:cnb_complete}, ${\mathcal{C}}'$ satisfies Definitions~\ref{def:con_bcc}, \ref{def:z_pf}, and~\ref{def:y_cl}.
%
We show that Algorithm~\ref{alg:cnb_syn} cannot remove all states of ${\mathcal{C}}'$. 
Consider the BCC ${\mathcal{C}}'$ obtained after each pruning iteration, and maintain the invariant component of ${\mathcal{C}}'$ that all states of ${\mathcal{C}}'$ are still contained in the current BCC. 
One can easily check that this invariant component holds initially.
%
Suppose that the invariant component holds before one pruning iteration. 
Since the Z-states within ${\mathcal{C}}'$ satisfy Definition~\ref{def:z_pf}, none of its Z-states can be removed as an infeasible Z-state. 
Considering the Y-states ${\mathcal{C}}'$ satisfy Definition~\ref{def:y_cl}, none of its Y-states can be removed as a non-cognitively-live Y-state. 
Meanwhile, since ${\mathcal{C}}'$ satisfies Definition~\ref{def:con_bcc}, none of its Y- or Z-states can be removed due to a consistency violation. 
Thus, the invariant component is preserved after this iteration.

Therefore, after the pruning procedure terminates, the remaining BCC still contains the accessible BCC representation ${\mathcal{C}}'$. 
In particular, the accessible component from $q^Y_0$ is nonempty and contains a cognitively marked Y-microstate. 
Finally, Algorithm~\ref{alg:cnb_syn} cannot return ``No solution''.
\end{proof}
%

\begin{figure}[ht!]
    \centering
    \includegraphics[width=0.625\linewidth]{Figs/controlled_fsa_overall.pdf}
    \caption{Controlled DFA $S / G$.}    % TODO 这个符号对不对
    \label{fig:ctrl_DFA}
\end{figure}

We give Example~\ref{ex:DFA_cnb} to illustrate that the closed-loop system satisfies the notion of cognitive non-blocking.
Meanwhile, we discuss the complexity of Algorithm~\ref{alg:cnb_syn}.

\begin{example}\label{ex:DFA_cnb}
    Let us consider the DFA illustrated in Fig.~\ref{fig:dfa_grid_env} (b) again, and the corresponding BCC illustrated in Fig.~\ref{fig:bcc_overall}.
    %
    After pruning the inconsistent states, the corresponding closed-loop system $S / G$ is illustrated in Fig.~\ref{fig:ctrl_DFA}.
    %
    In the resulting closed-loop system, the evolutions from the initial state $q_0$ are $q_0 \xrightarrow{u_1} q_1 \xrightarrow{o_2} q_5$ and $q_0 \xrightarrow{u_2} q_2 \xrightarrow{o_2} q_6$, and the corresponding observation is $o_2$.
    In the cognitive sense, after observing $o_2$, the supervisor can certify that the current DFA state belongs to the marked set
    \begin{equation}\label{eq:ex_cnb_reach_marked}
        \hat{\mathcal{E}}_S(o_2) = \{ q_5, q_6 \} \subseteq Q_m,
    \end{equation}
    although the exact state may still be unknown.
    %
    Moreover, from any reachable state after $q_5$ or $q_6$, the controlled system can evolve through a finite admissible continuation and return to $q_0$.
    The particular observable labels along this return segment are not essential for the cognitive non-blocking argument.
    After returning to $q_0$, the continuation for states of \eqref{eq:ex_cnb_reach_marked} can be applied again to reach the cognitively marked information state $\{ q_5, q_6 \}$.
    %
    Hence, for any sequence in $G$ corresponding to a cognitively live path in the BCC $\mathcal{C}$, there exists a continuation that reaches a cognitively marked state.
    Once the supervisor $S$ is induced from the BCC $\mathcal{C}$ as described in Section~\ref{subsec:supervisor_induction}, the closed-loop system $S/G$ inherits this property and satisfies cognitive non-blocking.
    
    Let us consider the motivating example again in Section~\ref{sec:intro}.
    The robotic system cannot certify its exact current state, but its state estimate can be updated after completing certain motion primitives.
    When it is at the initial state $q_0$, it should reject the motions corresponding to $o_3$, which may lead to $q_3$ or $q_4$.
    In this sense, it can then arrive at $q_5$ or $q_6$; moreover, when it observes the information corresponding to event $o_2$, it knows that it is in a safe location, so the transmission of critical information can be started.
\end{example}

\begin{proposition}
Consider a BCC $\mathcal{C}$ constructed over a DFA $G$ with the number of $|Q|$ states and $|\Sigma|$ events. 
The worst-case time complexity of Algorithm~\ref{alg:cnb_syn} is
$
    O(|Q|\cdot |\Sigma|\cdot 2^{2|Q|+2|\Sigma|}).
$
\end{proposition}

\begin{proof}
Algorithm~\ref{alg:cnb_syn} proceeds by pruning Y- and Z-states of the BCC.
In Lines 1--2, deadlock and livelock Z-microstates are identified by a BFS
search, which can be done in time
$
O(|\bm{\mathcal{Q}}|+|\delta_{\mathcal{C}}|)
$.
%
In Lines 3--6, each iteration prunes Y-states that are not cognitively live
or Z-states that violate consistency. Each iteration can be implemented by one
BFS over the remaining BCC, with cost
$
O(|\bm{\mathcal{Q}}|+|\delta_{\mathcal{C}}|)
$.
Since each nontrivial iteration prunes at least one Y-state or Z-state, the
number of iterations is at most $|Y|+|Z|$. Therefore, the total cost is
$
    O((|Y|+|Z|)(|\bm{\mathcal{Q}}|+|\delta_{\mathcal{C}}|)).
$
Line 7 performs one additional BFS, which is dominated by the above term.
Hence, the overall complexity is
$
    O((|Y|+|Z|)(|\bm{\mathcal{Q}}|+|\delta_{\mathcal{C}}|)).
$
%
By Proposition~\ref{prop:state_bcc}, we have
$
    |\bm{\mathcal{Q}}|
    \le
    |Q|(2^{|Q|}+2^{|Q|+|\Sigma|})
    =
    O(|Q|\cdot 2^{|Q|+|\Sigma|}),
$
and
$$
    |\delta_{\mathcal{C}}|
    \le
    |Q|2^{|Q|+|\Sigma|}(1+|\Sigma|)
    =
    O(|Q|\cdot |\Sigma|\cdot 2^{|Q|+|\Sigma|}).
$$
Therefore, based on
$
|Y|+|Z| = O(2^{|Q|+|\Sigma|}),
$
we have
$$
\begin{aligned}
    O((|Y|+|Z|)(|\bm{\mathcal{Q}}|+|\delta_{\mathcal{C}}|))
    &=
    O(
    2^{|Q|+|\Sigma|}
    \cdot
    |Q|\cdot |\Sigma|\cdot 2^{|Q|+|\Sigma|}
    ) \\
    &=
    O(|Q|\cdot |\Sigma|\cdot 2^{2|Q|+2|\Sigma|}).
\end{aligned}
$$
This completes the proof.
\end{proof}
%

As for the maximally permissiveness of the synthesized supervisor, the comparison of control patterns should be made at the same decision instant. Let
$
    \iota = y_0 z_0 y_1 z_1 \cdots y_n z_n \cdots
$
be an information-state sequence of BCC $\mathcal{C}$ generated by the synthesized supervisor \eqref{eq:bcc_sc:1}, where $\gamma_i$ is the control pattern made at Y-state $y_i$.
The maximally permissiveness requires that for any other $\gamma_i' \in H(y_i)$, we have $\gamma_i \not \subset \gamma_i'$.
%
In this sense, the maximally permissive partial-observation supervisor can be defined as
\begin{flalign}\label{eq:mps_cnb}
    S(\alpha) = \{ &\nu = \gamma_0 \gamma_1 \cdots \gamma_n \cdots \in \Gamma^* \mid \exists \rho_{\mathcal{C}} \in \mathcal{R}_{\mathcal{C}}, 
    \bm{s} = tr_{\mathcal{C}}(\rho_{\mathcal{C}}) :  \notag \\
    &\alpha = P(P_\Sigma(\bm{s})), \nu = P_\Gamma(\bm{s}), \iota = y_0 z_0 y_1 z_1 \cdots y_n   \notag \\
    & z_n \dots = R_X(\rho_{\mathcal{C}}),
    \forall \gamma_i' \in H(y_i), \gamma_i \not \subseteq \gamma_i'; \},
\end{flalign}
where $i \in \mathbb{N}$.

\begin{corollary}
Problem~\ref{prob:cnb} can be solved within finite steps by Algorithm~\ref{alg:cnb_syn}, and the maximally permissive supervisor can be obtained from \eqref{eq:mps_cnb}.
\end{corollary}


\section{Special case: marked tasks within the unobservable reaches}
\label{sec:special_z_marked}

%
% 这个系统要处理的话
% 其实方法是一样的
% 1 对于marked Z-state，需要当前Z-状态中的所有状态能够到达这个状态，或者退到下一个Y状态
% 2 Y状态不变
% 3 系统要consistent
% 

In this section, we discuss a special case of the proposed method, namely, systems in which marked tasks may be completed within unobservable reaches.
In the previous sections, the cognitively marked states are defined over Y-microstates, since Y-states represent information states immediately after observation updates.
However, in some systems, a marked state can be reached inside a Z-state before the next observable event occurs.
%
The main result of this article already handles the case where task completion is certified at marked Y-microstates. Therefore, we only briefly discuss the case of marked Z-microstates here.

To handle this case, we extend the target set of the BCC by including marked Z-microstates. 
For a Z-state $z$, define
$
    \bm{\mathcal{Q}}_{z,m} = \{(z,q)\in \bm{\mathcal{Q}}_z \mid q\in Q_m\}
$
as the set of marked Z-microstates in $z$. 
Meanwhile, the set of exit Z-microstates can be defined as
$$
    \bm{\mathcal{Q}}_{z,e}
    =
    \{ (z, q_e) \in \bm{\mathcal{Q}}_z \mid
    \exists \sigma_o \in (\Gamma(z) \cup \Sigma_{uc}) \cap \Sigma_o : \delta(q_e,\sigma_o)! \}.
$$
%
We say that a Z-state $z$ is a \emph{marked Z-state} if (i) for every microstate $(z,q)\in \bm{\mathcal{Q}}_z$, there exists a marked Z-microstate $(z,q_m)\in \bm{\mathcal{Q}}_{z,m}$ such that
$$
    (z,q) \leadsto_{\bm{\mathcal{Q}}_z} (z,q_m).
$$
(ii) for this marked Z-microstate $(z, q_m)$, there exists an exit Z-microstate $(z, q_e)\in \bm{\mathcal{Q}}_{z,e}$ and a path inside $\bm{\mathcal{Q}}_z$
$$
    \rho_{\mathcal{C}}
    =
    (z,q_0)(z,q_1)\cdots(z,q_\ell),
    \quad
    q_0=q_m, \, q_\ell=q_e,
$$
such that, for each $i=0,\ldots,\ell-1$, there exists $\sigma_{uo,i} \in \Sigma_{uo}$ satisfying
$
    \delta_{\mathcal{C}}((z,q_i),\sigma_{uo,i}) = (z,q_{i+1}),
$
and
$
    Post((z,q_i))\cap \bm{\mathcal{Q}}_z = \{ (z, q_{i+1}) \}.
$
Hence, after reaching the marked Z-microstate $(z,q_m)$, the subsequent internal evolution inside $z$ is uniquely determined until the exit Z-microstate $(z,q_e)$ is reached.
%
In other words, each microstate of a marked Z-state can reach a marked Z-microstate, after which the internal path to the exit is unique. 
Thus, the observed exit can certify task completion within the unobservable reach.

We use $Z_m$ to denote the set of marked Z-states.
Then, the target set for cognitive-liveness is extended from
$\bm{\mathcal{Q}}_m$ to
\begin{equation}
    \bm{\mathcal{Q}}_m^{+}
    =
    \bm{\mathcal{Q}}_m \cup \bigcup_{z\in Z_m}\bm{\mathcal{Q}}_{z,m}.
    \label{eq:extended_target_set}
\end{equation}
Accordingly, the definition of cognitively live Y-states in
\eqref{eq:bcc_y_cl} is extended as
\begin{equation}
\label{eq:bcc_y_cl:2}
    \forall q \in y,\;
    \exists (\bm{q}_t,q_t)\in \bm{\mathcal{Q}}_m^{+} :
    (\bm{q}_t,q_t) = \delta_{\mathcal{C}}((y,q),tr_{\mathcal{C}}(\rho_{\mathcal{C}}')),
\end{equation}
where the same BCC sequence $tr_{\mathcal{C}}(\rho_{\mathcal{C}}')$ is used for all $q\in y$, as in \eqref{eq:bcc_y_cl}.
Thus, a cognitively live Y-state is allowed to reach either a cognitively marked Y-microstate or a marked Z-microstate with an internal path to the next observable exit.
%
Based on these modifications, Algorithm~\ref{alg:cnb_syn} can still be applied to obtain a consistent BCC. 
The resulting BCC then induces a cognitive non-blocking supervisor for systems in which task completion can be certified either at marked Y-microstates or within marked Z-states.
In this way, Problem~\ref{prob:cnb} can be handled by replacing \eqref{eq:est_o} with \eqref{eq:esti:1}.

It should be noted that completion of a marked Z-state is certified only after the subsequent observable exit is observed. 
This is sufficient for the mathematical definition, since all compatible BCC paths must have visited a marked Z-microstate before the exit. 
However, in many practical applications, the unobservable tail inside the Z-state is not physically observed, and thus the exact occurrence of task completion may not be directly interpretable from the subsequent observation alone.

\section{Conclusion}\label{sec:con}

In this article, we proposed a new non-blocking property, called cognitive non-blocking, under which the observer can certify, from partial observations, that the system has reached a marked state.
Based on this notion, we developed a supervisor synthesis algorithm that constructs a controlled sublanguage guaranteed to be cognitive non-blocking. 
We also proved that the proposed synthesis algorithm terminates in finite steps. 
%
Future work will investigate how the notion of cognitive non-blocking can be extended to weighted and interval-weighted systems, where task completion may depend not only on reachability but also on accumulated
costs or weights.

\section*{References}\vspace{-1em}
\bibliographystyle{IEEEtran}
\bibliography{ref_abbr}


\begin{IEEEbiography}[{\includegraphics[width=1in,height=1.25in,clip,keepaspectratio]{a1.png}}]{First A. Author} (Fellow, IEEE) and all authors may include 
biographies. Biographies are
often not included in conference-related articles.
This author is an IEEE Fellow. The first paragraph
may contain a place and/or date of birth (list
place, then date). Next, the author’s educational
background is listed. The degrees should be listed
with type of degree in what field, which institution,
city, state, and country, and year the degree was
earned. The author’s major field of study should
be lower-cased.

The second paragraph uses the pronoun of the person (he or she) and
not the author’s last name. It lists military and work experience, including
summer and fellowship jobs. Job titles are capitalized. The current job must
have a location; previous positions may be listed without one. Information
concerning previous publications may be included. Try not to list more than
three books or published articles. The format for listing publishers of a book
within the biography is: title of book (publisher name, year) similar to a
reference. Current and previous research interests end the paragraph.

The third paragraph begins with the author’s title and last name (e.g.,
Dr. Smith, Prof. Jones, Mr. Kajor, Ms. Hunter). List any memberships in
professional societies other than the IEEE. Finally, list any awards and work
for IEEE committees and publications. If a photograph is provided, it should
be of good quality, and professional-looking.
\end{IEEEbiography}

\begin{IEEEbiographynophoto}{Second B. Author,} photograph and biography not available at the
time of publication.
\end{IEEEbiographynophoto}

\begin{IEEEbiographynophoto}{Third C. Author Jr.} (Member, IEEE), photograph and biography not available at the
time of publication.
\end{IEEEbiographynophoto}

\end{document}
